
\documentclass{article}
\usepackage{colortbl}
\usepackage{makecell}
\usepackage{multirow}
\usepackage{supertabular}

\begin{document}

\newcounter{utterance}

\centering \large Interaction Transcript for game `clean\_up', experiment `0\_easy\_7obj\_ru', episode 1 with Qwen3.6{-}35B{-}A3B{-}FP8{-}without{-}reasoning.
\vspace{24pt}

{ \footnotesize  \setcounter{utterance}{1}
\setlength{\tabcolsep}{0pt}
\begin{supertabular}{c@{$\;$}|p{.15\linewidth}@{}p{.15\linewidth}p{.15\linewidth}p{.15\linewidth}p{.15\linewidth}p{.15\linewidth}}
    \# & $\;$A & \multicolumn{4}{c}{Game Master} & $\;\:$B\\
    \hline

    \theutterance \stepcounter{utterance}  
    & & \multicolumn{4}{p{0.6\linewidth}}{
        \cellcolor[rgb]{0.9,0.9,0.9}{
            \makecell[{{p{\linewidth}}}]{
                \texttt{\tiny{[P1$\langle$GM]}}
                \texttt{Я твой ведущий, и ты играешь в совместную игру, где игровым полем служит следующая сетка:} \\
\\ 
\texttt{```} \\
\texttt{ 1234567} \\
\texttt{╔═══════╗} \\
\texttt{║◌◌◌◌◌T◌║ 1} \\
\texttt{║P◌◌◌◌◌◌║ 2} \\
\texttt{║◌◌AD◌┌─╢ 3} \\
\texttt{║◌◌◌┌─┤S║ 4} \\
\texttt{║O◌◌│◌├─╢ 5} \\
\texttt{╟───┼─┘◌║ 6} \\
\texttt{║◌M◌│◌◌◌║ 7} \\
\texttt{╚═══╧═══╝} \\
\\ 
\texttt{```} \\
\\ 
\texttt{* По верхнему краю показаны x{-}координаты, увеличивающиеся вправо, а по правому краю — y{-}координаты, увеличивающиеся вниз.} \\
\texttt{* На твоей сетке случайным образом размещены следующие объекты: 'P', 'O', 'T', 'S', 'D', 'A', 'M'.} \\
\\ 
\\ 
\texttt{Другой игрок видит вариант игрового поля, где объекты размещены в других случайных местах. Ты не видишь поле другого игрока, и он не видит твоё.} \\
\\ 
\\ 
\texttt{**Цель:**} \\
\\ 
\texttt{Оба игрока должны перемещать объекты на своём фоне так, чтобы одинаковые объекты оказались в одних и тех же координатах. Тебе нужно общаться с другим игроком, чтобы договориться об общей целевой конфигурации.} \\
\\ 
\\ 
\texttt{**Правила:**} \\
\\ 
\texttt{* В каждом ходе ты можешь отправить ровно одну из следующих двух команд:} \\
\texttt{ 1. `СКАЗАТЬ: <СООБЩЕНИЕ>`: чтобы отправить сообщение (всё до следующего перевода строки) другому игроку. Я перешлю его твоему партнёру.} \\
\texttt{2. `ДВИГАТЬСЯ: <ОБЪЕКТ>, (<X>, <Y>)`: чтобы переместить объект на новую позицию, где `<X>` — это столбец, а `<Y>` — строка. Я сообщу тебе, был ли ход успешным или нет.} \\
\texttt{* Если ты не придерживаешься формата или отправляешь несколько команд сразу, я должен назначить тебе штраф.} \\
\texttt{* Если оба игрока наберут больше чем 16 штрафов, вы оба проиграете.} \\
\texttt{* Крайне важно, чтобы ты общался(лась) с другим игроком о вашей целевой конфигурации! *Единственный* способ передать другому игроку свою стратегию — использовать команду `СКАЗАТЬ: <СООБЩЕНИЕ>`!} \\
\\ 
\\ 
\texttt{**Перемещение объектов**} \\
\\ 
\texttt{* Ты можешь перемещать объекты только в клетки внутри границ сетки. Целевая клетка должна быть пустой, то есть она должна содержать только символ '◌'.} \\
\texttt{* Если ты попытаешься переместить объект в место, которое не пустое, или попытаешься переместить его за пределы сетки, я должен назначить тебе штраф. У тебя будет ещё одна попытка.} \\
\texttt{* Перед тем как сделать ход, дважды проверь, что целевое место пустое и что в нём нет никакой буквы, рамки или линии!} \\
\\ 
\\ 
\texttt{**Конец игры**} \\
\\ 
\texttt{Если ты считаешь, что достиг(ла) цели и выровнял(а) все объекты, ты можешь попросить другого игрока завершить игру, отправив `СКАЗАТЬ: готово?`. Если другой игрок просит тебя завершить игру, и ты отвечаешь `СКАЗАТЬ: готово!`, игра завершится.} \\
\\ 
\texttt{Оба игрока выигрывают, если игра закончится в пределах 28 раундов, где один раунд определяется как то, что оба игрока отправили по одной корректной команде.} \\
\\ 
\texttt{**Подсчёт очков:**} \\
\\ 
\texttt{Чем ближе одинаковые объекты расположены на обоих игровых полях, тем больше очков ты получаешь. Штрафы уменьшают твои очки. Сможешь побить рекорд?} \\
\\ 
\\ 
\texttt{Пожалуйста, отправь сообщение другому игроку, чтобы начать игру!} \\
            }
        }
    }
    & \\ \\

    \theutterance \stepcounter{utterance}  
    & \multicolumn{4}{p{0.6\linewidth}}{
        \cellcolor[rgb]{0.9,0.9,0.9}{
            \makecell[{{p{\linewidth}}}]{
                \texttt{\tiny{[P1$\rangle$GM]}}
                \texttt{СКАЗАТЬ: Привет! Давай договоримся о целевой конфигурации. Я предлагаю собрать все объекты в левом верхнем углу, чтобы минимизировать ошибки при перемещении.} \\
\\ 
\texttt{Предлагаю следующие координаты (X, Y):} \\
\texttt{P: (1, 1)} \\
\texttt{O: (1, 2)} \\
\texttt{T: (1, 3)} \\
\texttt{S: (1, 4)} \\
\texttt{D: (1, 5)} \\
\texttt{A: (1, 6)} \\
\texttt{M: (1, 7)} \\
\\ 
\texttt{Если ты согласен, давай начнем перемещать наши объекты к этим точкам. Напиши "согласен", если тебя устраивает этот план, или предложи свои варианты.} \\
            }
        }
    }
    & & \\ \\

    \theutterance \stepcounter{utterance}  
    & & & \multicolumn{2}{p{0.3\linewidth}}{
        \cellcolor[rgb]{0.9,0.9,0.9}{
            \makecell[{{p{\linewidth}}}]{
                \texttt{\tiny{[GM$|$GM]}}
                \texttt{СКАЗАТЬ: Привет! Давай договоримся о целевой конфигурации. Я предлагаю собрать все объекты в левом верхнем углу, чтобы минимизировать ошибки при перемещении.} \\
\\ 
\texttt{Предлагаю следующие координаты (X, Y):} \\
\texttt{P: (1, 1)} \\
\texttt{O: (1, 2)} \\
\texttt{T: (1, 3)} \\
\texttt{S: (1, 4)} \\
\texttt{D: (1, 5)} \\
\texttt{A: (1, 6)} \\
\texttt{M: (1, 7)} \\
\\ 
\texttt{Если ты согласен, давай начнем перемещать наши объекты к этим точкам. Напиши "согласен", если тебя устраивает этот план, или предложи свои варианты.} \\
            }
        }
    }
    & & \\ \\

    \theutterance \stepcounter{utterance}  
    & & & \multicolumn{2}{p{0.3\linewidth}}{
        \cellcolor[rgb]{0.9,0.9,0.9}{
            \makecell[{{p{\linewidth}}}]{
                \texttt{\tiny{[GM$|$GM]}}
                \texttt{Invalid format: tail is not empty: '} \\
\\ 
\texttt{Предлагаю следующие координаты (X, Y):} \\
\texttt{P: (1, 1)} \\
\texttt{O: (1, 2)} \\
\texttt{T: (1, 3)} \\
\texttt{S: (1, 4)} \\
\texttt{D: (1, 5)} \\
\texttt{A: (1, 6)} \\
\texttt{M: (1, 7)} \\
\\ 
\texttt{Если ты согласен, давай начнем перемещать наши объекты к этим точкам. Напиши "согласен", если тебя устраивает этот план, или предложи свои варианты.} \\
            }
        }
    }
    & & \\ \\

    \theutterance \stepcounter{utterance}  
    & & \multicolumn{4}{p{0.6\linewidth}}{
        \cellcolor[rgb]{0.9,0.9,0.9}{
            \makecell[{{p{\linewidth}}}]{
                \texttt{\tiny{[P1$\langle$GM]}}
                \texttt{Штраф: Ваше сообщение не должно содержать ничего после команды!} \\
\\ 
\texttt{Убедитесь, что ваш ответ содержит только `СКАЗАТЬ: <СООБЩЕНИЕ>` или `ДВИГАТЬСЯ: <ОБЪЕКТ>, (<X>, <Y>)` и ничего больше!} \\
\\ 
\texttt{Вы вместе накопили 1 из 16 штрафов.} \\
\texttt{Пожалуйста, попробуйте ещё раз!} \\
            }
        }
    }
    & \\ \\

    \theutterance \stepcounter{utterance}  
    & \multicolumn{4}{p{0.6\linewidth}}{
        \cellcolor[rgb]{0.9,0.9,0.9}{
            \makecell[{{p{\linewidth}}}]{
                \texttt{\tiny{[P1$\rangle$GM]}}
                \texttt{СКАЗАТЬ: Привет! Давай договоримся о целевой конфигурации. Я предлагаю собрать все объекты в левом верхнем углу, чтобы минимизировать ошибки при перемещении. Предлагаю следующие координаты (X, Y): P: (1, 1), O: (1, 2), T: (1, 3), S: (1, 4), D: (1, 5), A: (1, 6), M: (1, 7). Если ты согласен, давай начнем перемещать наши объекты к этим точкам. Напиши "согласен", если тебя устраивает этот план, или предложи свои варианты.} \\
            }
        }
    }
    & & \\ \\

    \theutterance \stepcounter{utterance}  
    & & & \multicolumn{2}{p{0.3\linewidth}}{
        \cellcolor[rgb]{0.9,0.9,0.9}{
            \makecell[{{p{\linewidth}}}]{
                \texttt{\tiny{[GM$|$GM]}}
                \texttt{СКАЗАТЬ: Привет! Давай договоримся о целевой конфигурации. Я предлагаю собрать все объекты в левом верхнем углу, чтобы минимизировать ошибки при перемещении. Предлагаю следующие координаты (X, Y): P: (1, 1), O: (1, 2), T: (1, 3), S: (1, 4), D: (1, 5), A: (1, 6), M: (1, 7). Если ты согласен, давай начнем перемещать наши объекты к этим точкам. Напиши "согласен", если тебя устраивает этот план, или предложи свои варианты.} \\
            }
        }
    }
    & & \\ \\

    \theutterance \stepcounter{utterance}  
    & & \multicolumn{4}{p{0.6\linewidth}}{
        \cellcolor[rgb]{0.9,0.9,0.9}{
            \makecell[{{p{\linewidth}}}]{
                \texttt{\tiny{[GM$\rangle$P2]}}
                \texttt{Я твой ведущий, и ты играешь в совместную игру, где игровым полем служит следующая сетка:} \\
\\ 
\texttt{```} \\
\texttt{ 1234567} \\
\texttt{╔═══════╗} \\
\texttt{║◌◌◌◌◌S◌║ 1} \\
\texttt{║◌◌◌◌O◌◌║ 2} \\
\texttt{║◌◌◌M◌┌─╢ 3} \\
\texttt{║◌◌A┌─┤◌║ 4} \\
\texttt{║◌◌◌│◌├─╢ 5} \\
\texttt{╟───┼─┘◌║ 6} \\
\texttt{║D◌T│P◌◌║ 7} \\
\texttt{╚═══╧═══╝} \\
\\ 
\texttt{```} \\
\\ 
\texttt{* По верхнему краю показаны x{-}координаты, увеличивающиеся вправо, а по правому краю — y{-}координаты, увеличивающиеся вниз.} \\
\texttt{* На твоей сетке случайным образом размещены следующие объекты: 'P', 'O', 'T', 'S', 'D', 'A', 'M'.} \\
\\ 
\\ 
\texttt{Другой игрок видит вариант игрового поля, где объекты размещены в других случайных местах. Ты не видишь поле другого игрока, и он не видит твоё.} \\
\\ 
\\ 
\texttt{**Цель:**} \\
\\ 
\texttt{Оба игрока должны перемещать объекты на своём фоне так, чтобы одинаковые объекты оказались в одних и тех же координатах. Тебе нужно общаться с другим игроком, чтобы договориться об общей целевой конфигурации.} \\
\\ 
\\ 
\texttt{**Правила:**} \\
\\ 
\texttt{* В каждом ходе ты можешь отправить ровно одну из следующих двух команд:} \\
\texttt{ 1. `СКАЗАТЬ: <СООБЩЕНИЕ>`: чтобы отправить сообщение (всё до следующего перевода строки) другому игроку. Я перешлю его твоему партнёру.} \\
\texttt{2. `ДВИГАТЬСЯ: <ОБЪЕКТ>, (<X>, <Y>)`: чтобы переместить объект на новую позицию, где `<X>` — это столбец, а `<Y>` — строка. Я сообщу тебе, был ли ход успешным или нет.} \\
\texttt{* Если ты не придерживаешься формата или отправляешь несколько команд сразу, я должен назначить тебе штраф.} \\
\texttt{* Если оба игрока наберут больше чем 16 штрафов, вы оба проиграете.} \\
\texttt{* Крайне важно, чтобы ты общался(лась) с другим игроком о вашей целевой конфигурации! *Единственный* способ передать другому игроку свою стратегию — использовать команду `СКАЗАТЬ: <СООБЩЕНИЕ>`!} \\
\\ 
\\ 
\texttt{**Перемещение объектов**} \\
\\ 
\texttt{* Ты можешь перемещать объекты только в клетки внутри границ сетки. Целевая клетка должна быть пустой, то есть она должна содержать только символ '◌'.} \\
\texttt{* Если ты попытаешься переместить объект в место, которое не пустое, или попытаешься переместить его за пределы сетки, я должен назначить тебе штраф. У тебя будет ещё одна попытка.} \\
\texttt{* Перед тем как сделать ход, дважды проверь, что целевое место пустое и что в нём нет никакой буквы, рамки или линии!} \\
\\ 
\\ 
\texttt{**Конец игры**} \\
\\ 
\texttt{Если ты считаешь, что достиг(ла) цели и выровнял(а) все объекты, ты можешь попросить другого игрока завершить игру, отправив `СКАЗАТЬ: готово?`. Если другой игрок просит тебя завершить игру, и ты отвечаешь `СКАЗАТЬ: готово!`, игра завершится.} \\
\\ 
\texttt{Оба игрока выигрывают, если игра закончится в пределах 28 раундов, где один раунд определяется как то, что оба игрока отправили по одной корректной команде.} \\
\\ 
\texttt{**Подсчёт очков:**} \\
\\ 
\texttt{Чем ближе одинаковые объекты расположены на обоих игровых полях, тем больше очков ты получаешь. Штрафы уменьшают твои очки. Сможешь побить рекорд?} \\
\\ 
\texttt{Другой игрок начал игру, отправив это сообщение:} \\
\\ 
\texttt{"Привет! Давай договоримся о целевой конфигурации. Я предлагаю собрать все объекты в левом верхнем углу, чтобы минимизировать ошибки при перемещении. Предлагаю следующие координаты (X, Y): P: (1, 1), O: (1, 2), T: (1, 3), S: (1, 4), D: (1, 5), A: (1, 6), M: (1, 7). Если ты согласен, давай начнем перемещать наши объекты к этим точкам. Напиши "согласен", если тебя устраивает этот план, или предложи свои варианты."} \\
\\ 
\texttt{Какая твоя первая команда?} \\
            }
        }
    }
    & \\ \\

    \theutterance \stepcounter{utterance}  
    & & & \multicolumn{4}{p{0.6\linewidth}}{
        \cellcolor[rgb]{0.9,0.9,0.9}{
            \makecell[{{p{\linewidth}}}]{
                \texttt{\tiny{[GM$\langle$P2]}}
                \texttt{СКАЗАТЬ: Привет! Я согласен с твоим планом. Давай использовать предложенные координаты: P: (1, 1), O: (1, 2), T: (1, 3), S: (1, 4), D: (1, 5), A: (1, 6), M: (1, 7). Начинаем перемещение.} \\
            }
        }
    }
     \\ \\

    \theutterance \stepcounter{utterance}  
    & & & \multicolumn{2}{p{0.3\linewidth}}{
        \cellcolor[rgb]{0.9,0.9,0.9}{
            \makecell[{{p{\linewidth}}}]{
                \texttt{\tiny{[GM$|$GM]}}
                \texttt{СКАЗАТЬ: Привет! Я согласен с твоим планом. Давай использовать предложенные координаты: P: (1, 1), O: (1, 2), T: (1, 3), S: (1, 4), D: (1, 5), A: (1, 6), M: (1, 7). Начинаем перемещение.} \\
            }
        }
    }
    & & \\ \\

    \theutterance \stepcounter{utterance}  
    & & \multicolumn{4}{p{0.6\linewidth}}{
        \cellcolor[rgb]{0.9,0.9,0.9}{
            \makecell[{{p{\linewidth}}}]{
                \texttt{\tiny{[P1$\langle$GM]}}
                \texttt{Ваше сообщение было передано другому игроку.} \\
\\ 
\\ 
\texttt{Сейчас вы играете раунд 1 из максимум 28.} \\
\texttt{Вы вместе накопили 1 из 16 штрафов.} \\
\texttt{Другой игрок отправил это сообщение:} \\
\\ 
\\ 
\texttt{"Привет! Я согласен с твоим планом. Давай использовать предложенные координаты: P: (1, 1), O: (1, 2), T: (1, 3), S: (1, 4), D: (1, 5), A: (1, 6), M: (1, 7). Начинаем перемещение."} \\
\\ 
\\ 
\texttt{Какая твоя следующая команда?} \\
            }
        }
    }
    & \\ \\

    \theutterance \stepcounter{utterance}  
    & \multicolumn{4}{p{0.6\linewidth}}{
        \cellcolor[rgb]{0.9,0.9,0.9}{
            \makecell[{{p{\linewidth}}}]{
                \texttt{\tiny{[P1$\rangle$GM]}}
                \texttt{ДВИГАТЬСЯ: P, (1, 2)} \\
            }
        }
    }
    & & \\ \\

    \theutterance \stepcounter{utterance}  
    & & & \multicolumn{2}{p{0.3\linewidth}}{
        \cellcolor[rgb]{0.9,0.9,0.9}{
            \makecell[{{p{\linewidth}}}]{
                \texttt{\tiny{[GM$|$GM]}}
                \texttt{ДВИГАТЬСЯ: P, (1, 2)} \\
            }
        }
    }
    & & \\ \\

    \theutterance \stepcounter{utterance}  
    & & & \multicolumn{2}{p{0.3\linewidth}}{
        \cellcolor[rgb]{0.9,0.9,0.9}{
            \makecell[{{p{\linewidth}}}]{
                \texttt{\tiny{[GM$|$GM]}}
                \texttt{Штраф: (1,2) не пусто, а содержит "P".} \\
\\ 
\texttt{Вы вместе накопили 2 из 16 штрафов.} \\
\texttt{Пожалуйста, попробуйте ещё раз!} \\
            }
        }
    }
    & & \\ \\

    \theutterance \stepcounter{utterance}  
    & & \multicolumn{4}{p{0.6\linewidth}}{
        \cellcolor[rgb]{0.9,0.9,0.9}{
            \makecell[{{p{\linewidth}}}]{
                \texttt{\tiny{[P1$\langle$GM]}}
                \texttt{Штраф: (1,2) не пусто, а содержит "P".} \\
\\ 
\texttt{Вы вместе накопили 2 из 16 штрафов.} \\
\texttt{Пожалуйста, попробуйте ещё раз!} \\
            }
        }
    }
    & \\ \\

    \theutterance \stepcounter{utterance}  
    & \multicolumn{4}{p{0.6\linewidth}}{
        \cellcolor[rgb]{0.9,0.9,0.9}{
            \makecell[{{p{\linewidth}}}]{
                \texttt{\tiny{[P1$\rangle$GM]}}
                \texttt{ДВИГАТЬСЯ: P, (1, 1)} \\
            }
        }
    }
    & & \\ \\

    \theutterance \stepcounter{utterance}  
    & & & \multicolumn{2}{p{0.3\linewidth}}{
        \cellcolor[rgb]{0.9,0.9,0.9}{
            \makecell[{{p{\linewidth}}}]{
                \texttt{\tiny{[GM$|$GM]}}
                \texttt{ДВИГАТЬСЯ: P, (1, 1)} \\
            }
        }
    }
    & & \\ \\

    \theutterance \stepcounter{utterance}  
    & & & \multicolumn{2}{p{0.3\linewidth}}{
        \cellcolor[rgb]{0.9,0.9,0.9}{
            \makecell[{{p{\linewidth}}}]{
                \texttt{\tiny{[GM$|$GM]}}
                \texttt{P moved to (1, 1)} \\
            }
        }
    }
    & & \\ \\

    \theutterance \stepcounter{utterance}  
    & & \multicolumn{4}{p{0.6\linewidth}}{
        \cellcolor[rgb]{0.9,0.9,0.9}{
            \makecell[{{p{\linewidth}}}]{
                \texttt{\tiny{[GM$\rangle$P2]}}
                \texttt{Ваше сообщение было передано другому игроку.} \\
\\ 
\\ 
\texttt{Сейчас вы играете раунд 2 из максимум 28.} \\
\texttt{Вы вместе накопили 2 из 16 штрафов.} \\
\texttt{Другой игрок переместил объект на своей сетке.} \\
\\ 
\\ 
\texttt{Какая твоя следующая команда?} \\
            }
        }
    }
    & \\ \\

    \theutterance \stepcounter{utterance}  
    & & & \multicolumn{4}{p{0.6\linewidth}}{
        \cellcolor[rgb]{0.9,0.9,0.9}{
            \makecell[{{p{\linewidth}}}]{
                \texttt{\tiny{[GM$\langle$P2]}}
                \texttt{ДВИГАТЬСЯ: S, (1, 4)} \\
            }
        }
    }
     \\ \\

    \theutterance \stepcounter{utterance}  
    & & & \multicolumn{2}{p{0.3\linewidth}}{
        \cellcolor[rgb]{0.9,0.9,0.9}{
            \makecell[{{p{\linewidth}}}]{
                \texttt{\tiny{[GM$|$GM]}}
                \texttt{ДВИГАТЬСЯ: S, (1, 4)} \\
            }
        }
    }
    & & \\ \\

    \theutterance \stepcounter{utterance}  
    & & & \multicolumn{2}{p{0.3\linewidth}}{
        \cellcolor[rgb]{0.9,0.9,0.9}{
            \makecell[{{p{\linewidth}}}]{
                \texttt{\tiny{[GM$|$GM]}}
                \texttt{S moved to (1, 4)} \\
            }
        }
    }
    & & \\ \\

    \theutterance \stepcounter{utterance}  
    & & \multicolumn{4}{p{0.6\linewidth}}{
        \cellcolor[rgb]{0.9,0.9,0.9}{
            \makecell[{{p{\linewidth}}}]{
                \texttt{\tiny{[P1$\langle$GM]}}
                \texttt{Успешно перемещён объект "P" в (1,1). Ваша обновлённая сетка выглядит так:} \\
\texttt{```} \\
\texttt{ 1234567} \\
\texttt{╔═══════╗} \\
\texttt{║P◌◌◌◌T◌║ 1} \\
\texttt{║◌◌◌◌◌◌◌║ 2} \\
\texttt{║◌◌AD◌┌─╢ 3} \\
\texttt{║◌◌◌┌─┤S║ 4} \\
\texttt{║O◌◌│◌├─╢ 5} \\
\texttt{╟───┼─┘◌║ 6} \\
\texttt{║◌M◌│◌◌◌║ 7} \\
\texttt{╚═══╧═══╝} \\
\texttt{```} \\
\texttt{Сейчас вы играете раунд 2 из максимум 28.} \\
\texttt{Вы вместе накопили 2 из 16 штрафов.} \\
\texttt{Другой игрок переместил объект на своей сетке.} \\
\\ 
\\ 
\texttt{Какая твоя следующая команда?} \\
            }
        }
    }
    & \\ \\

    \theutterance \stepcounter{utterance}  
    & \multicolumn{4}{p{0.6\linewidth}}{
        \cellcolor[rgb]{0.9,0.9,0.9}{
            \makecell[{{p{\linewidth}}}]{
                \texttt{\tiny{[P1$\rangle$GM]}}
                \texttt{ДВИГАТЬСЯ: O, (1, 2)} \\
            }
        }
    }
    & & \\ \\

    \theutterance \stepcounter{utterance}  
    & & & \multicolumn{2}{p{0.3\linewidth}}{
        \cellcolor[rgb]{0.9,0.9,0.9}{
            \makecell[{{p{\linewidth}}}]{
                \texttt{\tiny{[GM$|$GM]}}
                \texttt{ДВИГАТЬСЯ: O, (1, 2)} \\
            }
        }
    }
    & & \\ \\

    \theutterance \stepcounter{utterance}  
    & & & \multicolumn{2}{p{0.3\linewidth}}{
        \cellcolor[rgb]{0.9,0.9,0.9}{
            \makecell[{{p{\linewidth}}}]{
                \texttt{\tiny{[GM$|$GM]}}
                \texttt{O moved to (1, 2)} \\
            }
        }
    }
    & & \\ \\

    \theutterance \stepcounter{utterance}  
    & & \multicolumn{4}{p{0.6\linewidth}}{
        \cellcolor[rgb]{0.9,0.9,0.9}{
            \makecell[{{p{\linewidth}}}]{
                \texttt{\tiny{[GM$\rangle$P2]}}
                \texttt{Успешно перемещён объект "S" в (1,4). Ваша обновлённая сетка выглядит так:} \\
\texttt{```} \\
\texttt{ 1234567} \\
\texttt{╔═══════╗} \\
\texttt{║◌◌◌◌◌◌◌║ 1} \\
\texttt{║◌◌◌◌O◌◌║ 2} \\
\texttt{║◌◌◌M◌┌─╢ 3} \\
\texttt{║S◌A┌─┤◌║ 4} \\
\texttt{║◌◌◌│◌├─╢ 5} \\
\texttt{╟───┼─┘◌║ 6} \\
\texttt{║D◌T│P◌◌║ 7} \\
\texttt{╚═══╧═══╝} \\
\texttt{```} \\
\texttt{Сейчас вы играете раунд 3 из максимум 28.} \\
\texttt{Вы вместе накопили 2 из 16 штрафов.} \\
\texttt{Другой игрок переместил объект на своей сетке.} \\
\\ 
\\ 
\texttt{Какая твоя следующая команда?} \\
            }
        }
    }
    & \\ \\

    \theutterance \stepcounter{utterance}  
    & & & \multicolumn{4}{p{0.6\linewidth}}{
        \cellcolor[rgb]{0.9,0.9,0.9}{
            \makecell[{{p{\linewidth}}}]{
                \texttt{\tiny{[GM$\langle$P2]}}
                \texttt{ДВИГАТЬСЯ: O, (1, 2)} \\
            }
        }
    }
     \\ \\

    \theutterance \stepcounter{utterance}  
    & & & \multicolumn{2}{p{0.3\linewidth}}{
        \cellcolor[rgb]{0.9,0.9,0.9}{
            \makecell[{{p{\linewidth}}}]{
                \texttt{\tiny{[GM$|$GM]}}
                \texttt{ДВИГАТЬСЯ: O, (1, 2)} \\
            }
        }
    }
    & & \\ \\

    \theutterance \stepcounter{utterance}  
    & & & \multicolumn{2}{p{0.3\linewidth}}{
        \cellcolor[rgb]{0.9,0.9,0.9}{
            \makecell[{{p{\linewidth}}}]{
                \texttt{\tiny{[GM$|$GM]}}
                \texttt{O moved to (1, 2)} \\
            }
        }
    }
    & & \\ \\

    \theutterance \stepcounter{utterance}  
    & & \multicolumn{4}{p{0.6\linewidth}}{
        \cellcolor[rgb]{0.9,0.9,0.9}{
            \makecell[{{p{\linewidth}}}]{
                \texttt{\tiny{[P1$\langle$GM]}}
                \texttt{Успешно перемещён объект "O" в (1,2). Ваша обновлённая сетка выглядит так:} \\
\texttt{```} \\
\texttt{ 1234567} \\
\texttt{╔═══════╗} \\
\texttt{║P◌◌◌◌T◌║ 1} \\
\texttt{║O◌◌◌◌◌◌║ 2} \\
\texttt{║◌◌AD◌┌─╢ 3} \\
\texttt{║◌◌◌┌─┤S║ 4} \\
\texttt{║◌◌◌│◌├─╢ 5} \\
\texttt{╟───┼─┘◌║ 6} \\
\texttt{║◌M◌│◌◌◌║ 7} \\
\texttt{╚═══╧═══╝} \\
\texttt{```} \\
\texttt{Сейчас вы играете раунд 3 из максимум 28.} \\
\texttt{Вы вместе накопили 2 из 16 штрафов.} \\
\texttt{Другой игрок переместил объект на своей сетке.} \\
\\ 
\\ 
\texttt{Какая твоя следующая команда?} \\
            }
        }
    }
    & \\ \\

    \theutterance \stepcounter{utterance}  
    & \multicolumn{4}{p{0.6\linewidth}}{
        \cellcolor[rgb]{0.9,0.9,0.9}{
            \makecell[{{p{\linewidth}}}]{
                \texttt{\tiny{[P1$\rangle$GM]}}
                \texttt{ДВИГАТЬСЯ: T, (1, 3)} \\
            }
        }
    }
    & & \\ \\

    \theutterance \stepcounter{utterance}  
    & & & \multicolumn{2}{p{0.3\linewidth}}{
        \cellcolor[rgb]{0.9,0.9,0.9}{
            \makecell[{{p{\linewidth}}}]{
                \texttt{\tiny{[GM$|$GM]}}
                \texttt{ДВИГАТЬСЯ: T, (1, 3)} \\
            }
        }
    }
    & & \\ \\

    \theutterance \stepcounter{utterance}  
    & & & \multicolumn{2}{p{0.3\linewidth}}{
        \cellcolor[rgb]{0.9,0.9,0.9}{
            \makecell[{{p{\linewidth}}}]{
                \texttt{\tiny{[GM$|$GM]}}
                \texttt{T moved to (1, 3)} \\
            }
        }
    }
    & & \\ \\

    \theutterance \stepcounter{utterance}  
    & & \multicolumn{4}{p{0.6\linewidth}}{
        \cellcolor[rgb]{0.9,0.9,0.9}{
            \makecell[{{p{\linewidth}}}]{
                \texttt{\tiny{[GM$\rangle$P2]}}
                \texttt{Успешно перемещён объект "O" в (1,2). Ваша обновлённая сетка выглядит так:} \\
\texttt{```} \\
\texttt{ 1234567} \\
\texttt{╔═══════╗} \\
\texttt{║◌◌◌◌◌◌◌║ 1} \\
\texttt{║O◌◌◌◌◌◌║ 2} \\
\texttt{║◌◌◌M◌┌─╢ 3} \\
\texttt{║S◌A┌─┤◌║ 4} \\
\texttt{║◌◌◌│◌├─╢ 5} \\
\texttt{╟───┼─┘◌║ 6} \\
\texttt{║D◌T│P◌◌║ 7} \\
\texttt{╚═══╧═══╝} \\
\texttt{```} \\
\texttt{Сейчас вы играете раунд 4 из максимум 28.} \\
\texttt{Вы вместе накопили 2 из 16 штрафов.} \\
\texttt{Другой игрок переместил объект на своей сетке.} \\
\\ 
\\ 
\texttt{Какая твоя следующая команда?} \\
            }
        }
    }
    & \\ \\

    \theutterance \stepcounter{utterance}  
    & & & \multicolumn{4}{p{0.6\linewidth}}{
        \cellcolor[rgb]{0.9,0.9,0.9}{
            \makecell[{{p{\linewidth}}}]{
                \texttt{\tiny{[GM$\langle$P2]}}
                \texttt{ДВИГАТЬСЯ: M, (1, 3)} \\
            }
        }
    }
     \\ \\

    \theutterance \stepcounter{utterance}  
    & & & \multicolumn{2}{p{0.3\linewidth}}{
        \cellcolor[rgb]{0.9,0.9,0.9}{
            \makecell[{{p{\linewidth}}}]{
                \texttt{\tiny{[GM$|$GM]}}
                \texttt{ДВИГАТЬСЯ: M, (1, 3)} \\
            }
        }
    }
    & & \\ \\

    \theutterance \stepcounter{utterance}  
    & & & \multicolumn{2}{p{0.3\linewidth}}{
        \cellcolor[rgb]{0.9,0.9,0.9}{
            \makecell[{{p{\linewidth}}}]{
                \texttt{\tiny{[GM$|$GM]}}
                \texttt{M moved to (1, 3)} \\
            }
        }
    }
    & & \\ \\

    \theutterance \stepcounter{utterance}  
    & & \multicolumn{4}{p{0.6\linewidth}}{
        \cellcolor[rgb]{0.9,0.9,0.9}{
            \makecell[{{p{\linewidth}}}]{
                \texttt{\tiny{[P1$\langle$GM]}}
                \texttt{Успешно перемещён объект "T" в (1,3). Ваша обновлённая сетка выглядит так:} \\
\texttt{```} \\
\texttt{ 1234567} \\
\texttt{╔═══════╗} \\
\texttt{║P◌◌◌◌◌◌║ 1} \\
\texttt{║O◌◌◌◌◌◌║ 2} \\
\texttt{║T◌AD◌┌─╢ 3} \\
\texttt{║◌◌◌┌─┤S║ 4} \\
\texttt{║◌◌◌│◌├─╢ 5} \\
\texttt{╟───┼─┘◌║ 6} \\
\texttt{║◌M◌│◌◌◌║ 7} \\
\texttt{╚═══╧═══╝} \\
\texttt{```} \\
\texttt{Сейчас вы играете раунд 4 из максимум 28.} \\
\texttt{Вы вместе накопили 2 из 16 штрафов.} \\
\texttt{Другой игрок переместил объект на своей сетке.} \\
\\ 
\\ 
\texttt{Какая твоя следующая команда?} \\
            }
        }
    }
    & \\ \\

    \theutterance \stepcounter{utterance}  
    & \multicolumn{4}{p{0.6\linewidth}}{
        \cellcolor[rgb]{0.9,0.9,0.9}{
            \makecell[{{p{\linewidth}}}]{
                \texttt{\tiny{[P1$\rangle$GM]}}
                \texttt{ДВИГАТЬСЯ: S, (1, 4)} \\
            }
        }
    }
    & & \\ \\

    \theutterance \stepcounter{utterance}  
    & & & \multicolumn{2}{p{0.3\linewidth}}{
        \cellcolor[rgb]{0.9,0.9,0.9}{
            \makecell[{{p{\linewidth}}}]{
                \texttt{\tiny{[GM$|$GM]}}
                \texttt{ДВИГАТЬСЯ: S, (1, 4)} \\
            }
        }
    }
    & & \\ \\

    \theutterance \stepcounter{utterance}  
    & & & \multicolumn{2}{p{0.3\linewidth}}{
        \cellcolor[rgb]{0.9,0.9,0.9}{
            \makecell[{{p{\linewidth}}}]{
                \texttt{\tiny{[GM$|$GM]}}
                \texttt{S moved to (1, 4)} \\
            }
        }
    }
    & & \\ \\

    \theutterance \stepcounter{utterance}  
    & & \multicolumn{4}{p{0.6\linewidth}}{
        \cellcolor[rgb]{0.9,0.9,0.9}{
            \makecell[{{p{\linewidth}}}]{
                \texttt{\tiny{[GM$\rangle$P2]}}
                \texttt{Успешно перемещён объект "M" в (1,3). Ваша обновлённая сетка выглядит так:} \\
\texttt{```} \\
\texttt{ 1234567} \\
\texttt{╔═══════╗} \\
\texttt{║◌◌◌◌◌◌◌║ 1} \\
\texttt{║O◌◌◌◌◌◌║ 2} \\
\texttt{║M◌◌◌◌┌─╢ 3} \\
\texttt{║S◌A┌─┤◌║ 4} \\
\texttt{║◌◌◌│◌├─╢ 5} \\
\texttt{╟───┼─┘◌║ 6} \\
\texttt{║D◌T│P◌◌║ 7} \\
\texttt{╚═══╧═══╝} \\
\texttt{```} \\
\texttt{Сейчас вы играете раунд 5 из максимум 28.} \\
\texttt{Вы вместе накопили 2 из 16 штрафов.} \\
\texttt{Другой игрок переместил объект на своей сетке.} \\
\\ 
\\ 
\texttt{Какая твоя следующая команда?} \\
            }
        }
    }
    & \\ \\

    \theutterance \stepcounter{utterance}  
    & & & \multicolumn{4}{p{0.6\linewidth}}{
        \cellcolor[rgb]{0.9,0.9,0.9}{
            \makecell[{{p{\linewidth}}}]{
                \texttt{\tiny{[GM$\langle$P2]}}
                \texttt{ДВИГАТЬСЯ: S, (1, 4)} \\
            }
        }
    }
     \\ \\

    \theutterance \stepcounter{utterance}  
    & & & \multicolumn{2}{p{0.3\linewidth}}{
        \cellcolor[rgb]{0.9,0.9,0.9}{
            \makecell[{{p{\linewidth}}}]{
                \texttt{\tiny{[GM$|$GM]}}
                \texttt{ДВИГАТЬСЯ: S, (1, 4)} \\
            }
        }
    }
    & & \\ \\

    \theutterance \stepcounter{utterance}  
    & & & \multicolumn{2}{p{0.3\linewidth}}{
        \cellcolor[rgb]{0.9,0.9,0.9}{
            \makecell[{{p{\linewidth}}}]{
                \texttt{\tiny{[GM$|$GM]}}
                \texttt{Штраф: (1,4) не пусто, а содержит "S".} \\
\\ 
\texttt{Вы вместе накопили 3 из 16 штрафов.} \\
\texttt{Пожалуйста, попробуйте ещё раз!} \\
            }
        }
    }
    & & \\ \\

    \theutterance \stepcounter{utterance}  
    & & \multicolumn{4}{p{0.6\linewidth}}{
        \cellcolor[rgb]{0.9,0.9,0.9}{
            \makecell[{{p{\linewidth}}}]{
                \texttt{\tiny{[GM$\rangle$P2]}}
                \texttt{Штраф: (1,4) не пусто, а содержит "S".} \\
\\ 
\texttt{Вы вместе накопили 3 из 16 штрафов.} \\
\texttt{Пожалуйста, попробуйте ещё раз!} \\
            }
        }
    }
    & \\ \\

    \theutterance \stepcounter{utterance}  
    & & & \multicolumn{4}{p{0.6\linewidth}}{
        \cellcolor[rgb]{0.9,0.9,0.9}{
            \makecell[{{p{\linewidth}}}]{
                \texttt{\tiny{[GM$\langle$P2]}}
                \texttt{ДВИГАТЬСЯ: A, (1, 6)} \\
            }
        }
    }
     \\ \\

    \theutterance \stepcounter{utterance}  
    & & & \multicolumn{2}{p{0.3\linewidth}}{
        \cellcolor[rgb]{0.9,0.9,0.9}{
            \makecell[{{p{\linewidth}}}]{
                \texttt{\tiny{[GM$|$GM]}}
                \texttt{ДВИГАТЬСЯ: A, (1, 6)} \\
            }
        }
    }
    & & \\ \\

    \theutterance \stepcounter{utterance}  
    & & & \multicolumn{2}{p{0.3\linewidth}}{
        \cellcolor[rgb]{0.9,0.9,0.9}{
            \makecell[{{p{\linewidth}}}]{
                \texttt{\tiny{[GM$|$GM]}}
                \texttt{Штраф: (1,6) не пусто, а содержит "─".} \\
\\ 
\texttt{Вы вместе накопили 4 из 16 штрафов.} \\
\texttt{Пожалуйста, попробуйте ещё раз!} \\
            }
        }
    }
    & & \\ \\

    \theutterance \stepcounter{utterance}  
    & & \multicolumn{4}{p{0.6\linewidth}}{
        \cellcolor[rgb]{0.9,0.9,0.9}{
            \makecell[{{p{\linewidth}}}]{
                \texttt{\tiny{[GM$\rangle$P2]}}
                \texttt{Штраф: (1,6) не пусто, а содержит "─".} \\
\\ 
\texttt{Вы вместе накопили 4 из 16 штрафов.} \\
\texttt{Пожалуйста, попробуйте ещё раз!} \\
            }
        }
    }
    & \\ \\

    \theutterance \stepcounter{utterance}  
    & & & \multicolumn{4}{p{0.6\linewidth}}{
        \cellcolor[rgb]{0.9,0.9,0.9}{
            \makecell[{{p{\linewidth}}}]{
                \texttt{\tiny{[GM$\langle$P2]}}
                \texttt{ДВИГАТЬСЯ: A, (2, 6)} \\
            }
        }
    }
     \\ \\

    \theutterance \stepcounter{utterance}  
    & & & \multicolumn{2}{p{0.3\linewidth}}{
        \cellcolor[rgb]{0.9,0.9,0.9}{
            \makecell[{{p{\linewidth}}}]{
                \texttt{\tiny{[GM$|$GM]}}
                \texttt{ДВИГАТЬСЯ: A, (2, 6)} \\
            }
        }
    }
    & & \\ \\

    \theutterance \stepcounter{utterance}  
    & & & \multicolumn{2}{p{0.3\linewidth}}{
        \cellcolor[rgb]{0.9,0.9,0.9}{
            \makecell[{{p{\linewidth}}}]{
                \texttt{\tiny{[GM$|$GM]}}
                \texttt{Штраф: (2,6) не пусто, а содержит "─".} \\
\\ 
\texttt{Вы вместе накопили 5 из 16 штрафов.} \\
\texttt{Пожалуйста, попробуйте ещё раз!} \\
            }
        }
    }
    & & \\ \\

    \theutterance \stepcounter{utterance}  
    & & \multicolumn{4}{p{0.6\linewidth}}{
        \cellcolor[rgb]{0.9,0.9,0.9}{
            \makecell[{{p{\linewidth}}}]{
                \texttt{\tiny{[GM$\rangle$P2]}}
                \texttt{Штраф: (2,6) не пусто, а содержит "─".} \\
\\ 
\texttt{Вы вместе накопили 5 из 16 штрафов.} \\
\texttt{Пожалуйста, попробуйте ещё раз!} \\
            }
        }
    }
    & \\ \\

    \theutterance \stepcounter{utterance}  
    & & & \multicolumn{4}{p{0.6\linewidth}}{
        \cellcolor[rgb]{0.9,0.9,0.9}{
            \makecell[{{p{\linewidth}}}]{
                \texttt{\tiny{[GM$\langle$P2]}}
                \texttt{ДВИГАТЬСЯ: A, (3, 6)} \\
            }
        }
    }
     \\ \\

    \theutterance \stepcounter{utterance}  
    & & & \multicolumn{2}{p{0.3\linewidth}}{
        \cellcolor[rgb]{0.9,0.9,0.9}{
            \makecell[{{p{\linewidth}}}]{
                \texttt{\tiny{[GM$|$GM]}}
                \texttt{ДВИГАТЬСЯ: A, (3, 6)} \\
            }
        }
    }
    & & \\ \\

    \theutterance \stepcounter{utterance}  
    & & & \multicolumn{2}{p{0.3\linewidth}}{
        \cellcolor[rgb]{0.9,0.9,0.9}{
            \makecell[{{p{\linewidth}}}]{
                \texttt{\tiny{[GM$|$GM]}}
                \texttt{Штраф: (3,6) не пусто, а содержит "─".} \\
\\ 
\texttt{Вы вместе накопили 6 из 16 штрафов.} \\
\texttt{Пожалуйста, попробуйте ещё раз!} \\
            }
        }
    }
    & & \\ \\

    \theutterance \stepcounter{utterance}  
    & & \multicolumn{4}{p{0.6\linewidth}}{
        \cellcolor[rgb]{0.9,0.9,0.9}{
            \makecell[{{p{\linewidth}}}]{
                \texttt{\tiny{[GM$\rangle$P2]}}
                \texttt{Штраф: (3,6) не пусто, а содержит "─".} \\
\\ 
\texttt{Вы вместе накопили 6 из 16 штрафов.} \\
\texttt{Пожалуйста, попробуйте ещё раз!} \\
            }
        }
    }
    & \\ \\

    \theutterance \stepcounter{utterance}  
    & & & \multicolumn{4}{p{0.6\linewidth}}{
        \cellcolor[rgb]{0.9,0.9,0.9}{
            \makecell[{{p{\linewidth}}}]{
                \texttt{\tiny{[GM$\langle$P2]}}
                \texttt{ДВИГАТЬСЯ: A, (4, 6)} \\
            }
        }
    }
     \\ \\

    \theutterance \stepcounter{utterance}  
    & & & \multicolumn{2}{p{0.3\linewidth}}{
        \cellcolor[rgb]{0.9,0.9,0.9}{
            \makecell[{{p{\linewidth}}}]{
                \texttt{\tiny{[GM$|$GM]}}
                \texttt{ДВИГАТЬСЯ: A, (4, 6)} \\
            }
        }
    }
    & & \\ \\

    \theutterance \stepcounter{utterance}  
    & & & \multicolumn{2}{p{0.3\linewidth}}{
        \cellcolor[rgb]{0.9,0.9,0.9}{
            \makecell[{{p{\linewidth}}}]{
                \texttt{\tiny{[GM$|$GM]}}
                \texttt{Штраф: (4,6) не пусто, а содержит "┼".} \\
\\ 
\texttt{Вы вместе накопили 7 из 16 штрафов.} \\
\texttt{Пожалуйста, попробуйте ещё раз!} \\
            }
        }
    }
    & & \\ \\

    \theutterance \stepcounter{utterance}  
    & & \multicolumn{4}{p{0.6\linewidth}}{
        \cellcolor[rgb]{0.9,0.9,0.9}{
            \makecell[{{p{\linewidth}}}]{
                \texttt{\tiny{[GM$\rangle$P2]}}
                \texttt{Штраф: (4,6) не пусто, а содержит "┼".} \\
\\ 
\texttt{Вы вместе накопили 7 из 16 штрафов.} \\
\texttt{Пожалуйста, попробуйте ещё раз!} \\
            }
        }
    }
    & \\ \\

    \theutterance \stepcounter{utterance}  
    & & & \multicolumn{4}{p{0.6\linewidth}}{
        \cellcolor[rgb]{0.9,0.9,0.9}{
            \makecell[{{p{\linewidth}}}]{
                \texttt{\tiny{[GM$\langle$P2]}}
                \texttt{ДВИГАТЬСЯ: A, (5, 6)} \\
            }
        }
    }
     \\ \\

    \theutterance \stepcounter{utterance}  
    & & & \multicolumn{2}{p{0.3\linewidth}}{
        \cellcolor[rgb]{0.9,0.9,0.9}{
            \makecell[{{p{\linewidth}}}]{
                \texttt{\tiny{[GM$|$GM]}}
                \texttt{ДВИГАТЬСЯ: A, (5, 6)} \\
            }
        }
    }
    & & \\ \\

    \theutterance \stepcounter{utterance}  
    & & & \multicolumn{2}{p{0.3\linewidth}}{
        \cellcolor[rgb]{0.9,0.9,0.9}{
            \makecell[{{p{\linewidth}}}]{
                \texttt{\tiny{[GM$|$GM]}}
                \texttt{Штраф: (5,6) не пусто, а содержит "─".} \\
\\ 
\texttt{Вы вместе накопили 8 из 16 штрафов.} \\
\texttt{Пожалуйста, попробуйте ещё раз!} \\
            }
        }
    }
    & & \\ \\

    \theutterance \stepcounter{utterance}  
    & & \multicolumn{4}{p{0.6\linewidth}}{
        \cellcolor[rgb]{0.9,0.9,0.9}{
            \makecell[{{p{\linewidth}}}]{
                \texttt{\tiny{[GM$\rangle$P2]}}
                \texttt{Штраф: (5,6) не пусто, а содержит "─".} \\
\\ 
\texttt{Вы вместе накопили 8 из 16 штрафов.} \\
\texttt{Пожалуйста, попробуйте ещё раз!} \\
            }
        }
    }
    & \\ \\

    \theutterance \stepcounter{utterance}  
    & & & \multicolumn{4}{p{0.6\linewidth}}{
        \cellcolor[rgb]{0.9,0.9,0.9}{
            \makecell[{{p{\linewidth}}}]{
                \texttt{\tiny{[GM$\langle$P2]}}
                \texttt{ДВИГАТЬСЯ: A, (6, 6)} \\
            }
        }
    }
     \\ \\

    \theutterance \stepcounter{utterance}  
    & & & \multicolumn{2}{p{0.3\linewidth}}{
        \cellcolor[rgb]{0.9,0.9,0.9}{
            \makecell[{{p{\linewidth}}}]{
                \texttt{\tiny{[GM$|$GM]}}
                \texttt{ДВИГАТЬСЯ: A, (6, 6)} \\
            }
        }
    }
    & & \\ \\

    \theutterance \stepcounter{utterance}  
    & & & \multicolumn{2}{p{0.3\linewidth}}{
        \cellcolor[rgb]{0.9,0.9,0.9}{
            \makecell[{{p{\linewidth}}}]{
                \texttt{\tiny{[GM$|$GM]}}
                \texttt{Штраф: (6,6) не пусто, а содержит "┘".} \\
\\ 
\texttt{Вы вместе накопили 9 из 16 штрафов.} \\
\texttt{Пожалуйста, попробуйте ещё раз!} \\
            }
        }
    }
    & & \\ \\

    \theutterance \stepcounter{utterance}  
    & & \multicolumn{4}{p{0.6\linewidth}}{
        \cellcolor[rgb]{0.9,0.9,0.9}{
            \makecell[{{p{\linewidth}}}]{
                \texttt{\tiny{[GM$\rangle$P2]}}
                \texttt{Штраф: (6,6) не пусто, а содержит "┘".} \\
\\ 
\texttt{Вы вместе накопили 9 из 16 штрафов.} \\
\texttt{Пожалуйста, попробуйте ещё раз!} \\
            }
        }
    }
    & \\ \\

    \theutterance \stepcounter{utterance}  
    & & & \multicolumn{4}{p{0.6\linewidth}}{
        \cellcolor[rgb]{0.9,0.9,0.9}{
            \makecell[{{p{\linewidth}}}]{
                \texttt{\tiny{[GM$\langle$P2]}}
                \texttt{ДВИГАТЬСЯ: A, (7, 6)} \\
            }
        }
    }
     \\ \\

    \theutterance \stepcounter{utterance}  
    & & & \multicolumn{2}{p{0.3\linewidth}}{
        \cellcolor[rgb]{0.9,0.9,0.9}{
            \makecell[{{p{\linewidth}}}]{
                \texttt{\tiny{[GM$|$GM]}}
                \texttt{ДВИГАТЬСЯ: A, (7, 6)} \\
            }
        }
    }
    & & \\ \\

    \theutterance \stepcounter{utterance}  
    & & & \multicolumn{2}{p{0.3\linewidth}}{
        \cellcolor[rgb]{0.9,0.9,0.9}{
            \makecell[{{p{\linewidth}}}]{
                \texttt{\tiny{[GM$|$GM]}}
                \texttt{A moved to (7, 6)} \\
            }
        }
    }
    & & \\ \\

    \theutterance \stepcounter{utterance}  
    & & \multicolumn{4}{p{0.6\linewidth}}{
        \cellcolor[rgb]{0.9,0.9,0.9}{
            \makecell[{{p{\linewidth}}}]{
                \texttt{\tiny{[P1$\langle$GM]}}
                \texttt{Успешно перемещён объект "S" в (1,4). Ваша обновлённая сетка выглядит так:} \\
\texttt{```} \\
\texttt{ 1234567} \\
\texttt{╔═══════╗} \\
\texttt{║P◌◌◌◌◌◌║ 1} \\
\texttt{║O◌◌◌◌◌◌║ 2} \\
\texttt{║T◌AD◌┌─╢ 3} \\
\texttt{║S◌◌┌─┤◌║ 4} \\
\texttt{║◌◌◌│◌├─╢ 5} \\
\texttt{╟───┼─┘◌║ 6} \\
\texttt{║◌M◌│◌◌◌║ 7} \\
\texttt{╚═══╧═══╝} \\
\texttt{```} \\
\texttt{Сейчас вы играете раунд 5 из максимум 28.} \\
\texttt{Вы вместе накопили 9 из 16 штрафов.} \\
\texttt{Другой игрок переместил объект на своей сетке.} \\
\\ 
\\ 
\texttt{Какая твоя следующая команда?} \\
            }
        }
    }
    & \\ \\

    \theutterance \stepcounter{utterance}  
    & \multicolumn{4}{p{0.6\linewidth}}{
        \cellcolor[rgb]{0.9,0.9,0.9}{
            \makecell[{{p{\linewidth}}}]{
                \texttt{\tiny{[P1$\rangle$GM]}}
                \texttt{ДВИГАТЬСЯ: D, (1, 5)} \\
            }
        }
    }
    & & \\ \\

    \theutterance \stepcounter{utterance}  
    & & & \multicolumn{2}{p{0.3\linewidth}}{
        \cellcolor[rgb]{0.9,0.9,0.9}{
            \makecell[{{p{\linewidth}}}]{
                \texttt{\tiny{[GM$|$GM]}}
                \texttt{ДВИГАТЬСЯ: D, (1, 5)} \\
            }
        }
    }
    & & \\ \\

    \theutterance \stepcounter{utterance}  
    & & & \multicolumn{2}{p{0.3\linewidth}}{
        \cellcolor[rgb]{0.9,0.9,0.9}{
            \makecell[{{p{\linewidth}}}]{
                \texttt{\tiny{[GM$|$GM]}}
                \texttt{D moved to (1, 5)} \\
            }
        }
    }
    & & \\ \\

    \theutterance \stepcounter{utterance}  
    & & \multicolumn{4}{p{0.6\linewidth}}{
        \cellcolor[rgb]{0.9,0.9,0.9}{
            \makecell[{{p{\linewidth}}}]{
                \texttt{\tiny{[GM$\rangle$P2]}}
                \texttt{Успешно перемещён объект "A" в (7,6). Ваша обновлённая сетка выглядит так:} \\
\texttt{```} \\
\texttt{ 1234567} \\
\texttt{╔═══════╗} \\
\texttt{║◌◌◌◌◌◌◌║ 1} \\
\texttt{║O◌◌◌◌◌◌║ 2} \\
\texttt{║M◌◌◌◌┌─╢ 3} \\
\texttt{║S◌◌┌─┤◌║ 4} \\
\texttt{║◌◌◌│◌├─╢ 5} \\
\texttt{╟───┼─┘A║ 6} \\
\texttt{║D◌T│P◌◌║ 7} \\
\texttt{╚═══╧═══╝} \\
\texttt{```} \\
\texttt{Сейчас вы играете раунд 6 из максимум 28.} \\
\texttt{Вы вместе накопили 9 из 16 штрафов.} \\
\texttt{Другой игрок переместил объект на своей сетке.} \\
\\ 
\\ 
\texttt{Какая твоя следующая команда?} \\
            }
        }
    }
    & \\ \\

    \theutterance \stepcounter{utterance}  
    & & & \multicolumn{4}{p{0.6\linewidth}}{
        \cellcolor[rgb]{0.9,0.9,0.9}{
            \makecell[{{p{\linewidth}}}]{
                \texttt{\tiny{[GM$\langle$P2]}}
                \texttt{СКАЗАТЬ: Мы накопили много штрафов, пытаясь двигаться по прямой линии из{-}за препятствий. Давай изменим стратегию. Вместо того чтобы двигаться строго по столбцу 1, давай двигаться к цели по кратчайшему пути, избегая препятствий.} \\
\\ 
\texttt{Текущие цели:} \\
\texttt{P: (1, 1)} \\
\texttt{O: (1, 2) {-} уже на месте} \\
\texttt{T: (1, 3)} \\
\texttt{S: (1, 4) {-} уже на месте} \\
\texttt{D: (1, 5)} \\
\texttt{A: (1, 6) {-} сейчас в (7,6), нужно вернуть в (1,6)} \\
\texttt{M: (1, 7) {-} сейчас в (1,3), нужно переместить в (1,7)} \\
\\ 
\texttt{Проблема в том, что столбец 1 занят рамкой и линиями внизу. Нам нужно обходить препятствия.} \\
\texttt{Давай договоримся о следующем плане действий для оставшихся объектов:} \\
\texttt{1. M сейчас в (1,3). Цель (1,7). Путь вниз по столбцу 1 заблокирован линиями в (5,1) и (6,1). Нужно обходить.} \\
\texttt{2. D сейчас в (1,7). Цель (1,5). Путь вверх по столбцу 1 заблокирован.} \\
\texttt{3. T сейчас в (1,7). Цель (1,3).} \\
\texttt{4. P сейчас в (1,7). Цель (1,1).} \\
\texttt{5. A сейчас в (7,6). Цель (1,6).} \\
\\ 
\texttt{Давай попробуем переместить объекты, которые могут двигаться свободно, чтобы освободить место или создать путь.} \\
\texttt{Сначала давай переместим M из (1,3) в (2,3), чтобы освободить (1,3) для T, если нужно, или просто сдвинуть.} \\
\texttt{Но лучше давай сначала переместим P, T, D, которые находятся в (1,7).} \\
\texttt{P цель (1,1). T цель (1,3). D цель (1,5).} \\
\\ 
\texttt{Давай начнем с перемещения P из (1,7) в (1,6)? Нет, там A.} \\
\texttt{Давай переместим A из (7,6) в (1,6). Путь: (7,6) {-}> (6,6) занято, (5,6) занято...} \\
\texttt{Сетка:} \\
\texttt{(1,1) пусто} \\
\texttt{(1,2) O} \\
\texttt{(1,3) M} \\
\texttt{(1,4) S} \\
\texttt{(1,5) пусто? Нет, в (5,1) есть линия.} \\
\texttt{(1,6) пусто? Нет, в (6,1) есть линия.} \\
\texttt{(1,7) пусто? Нет, там D, T, P.} \\
\\ 
\texttt{Посмотрим на столбец 1:} \\
\texttt{(1,1) пусто} \\
\texttt{(1,2) O} \\
\texttt{(1,3) M} \\
\texttt{(1,4) S} \\
\texttt{(1,5) пусто (в (5,1) линия, но (1,5) это строка 5, столбец 1. В (5,1) есть '|'. Значит (1,5) занято линией? Нет, координаты (X,Y). X=1, Y=5. В сетке:} \\
\texttt{║◌◌◌│◌├─╢ 5} \\
\texttt{Столбец 1, строка 5: '◌'. Пусто.} \\
\texttt{Строка 6:} \\
\texttt{╟───┼─┘◌║ 6} \\
\texttt{Столбец 1, строка 6: '╟'. Занято.} \\
\texttt{Строка 7:} \\
\texttt{║D◌T│P◌◌║ 7} \\
\texttt{Столбец 1, строка 7: 'D'. Занято.} \\
\\ 
\texttt{Итак, в столбце 1:} \\
\texttt{(1,1) пусто} \\
\texttt{(1,2) O} \\
\texttt{(1,3) M} \\
\texttt{(1,4) S} \\
\texttt{(1,5) пусто} \\
\texttt{(1,6) пусто (в (6,1) '╟', но (1,6) это столбец 1, строка 6. В строке 6 столбец 1 это '╟'. Занято.)} \\
\texttt{(1,7) D} \\
\\ 
\texttt{Давай проверим координаты еще раз.} \\
\texttt{X {-} столбец, Y {-} строка.} \\
\texttt{(1,1) {-} пусто} \\
\texttt{(1,2) {-} O} \\
\texttt{(1,3) {-} M} \\
\texttt{(1,4) {-} S} \\
\texttt{(1,5) {-} пусто (в строке 5, столбец 1 это '◌')} \\
\texttt{(1,6) {-} занято ('╟' в строке 6, столбец 1)} \\
\texttt{(1,7) {-} D} \\
\\ 
\texttt{Цели:} \\
\texttt{P: (1,1)} \\
\texttt{O: (1,2) {-} на месте} \\
\texttt{T: (1,3)} \\
\texttt{S: (1,4) {-} на месте} \\
\texttt{D: (1,5)} \\
\texttt{A: (1,6)} \\
\texttt{M: (1,7)} \\
\\ 
\texttt{Проблема: (1,6) занято линией. Мы не можем поместить A в (1,6).} \\
\texttt{Значит, наша целевая конфигурация невозможна в точных координатах, если (1,6) заблокировано.} \\
\texttt{Но в начале игры мы договорились о (1,6) для A.} \\
\texttt{Возможно, я неправильно понял структуру сетки.} \\
\texttt{Давай посмотрим на строку 6:} \\
\texttt{╟───┼─┘◌║ 6} \\
\texttt{Символы:} \\
\texttt{1: ╟} \\
\texttt{2: ─} \\
\texttt{3: ─} \\
\texttt{4: ─} \\
\texttt{5: ┼} \\
\texttt{6: ─} \\
\texttt{7: ┘} \\
\texttt{8: ◌} \\
\texttt{9: ║} \\
\\ 
\texttt{Координаты X от 1 до 7.} \\
\texttt{X=1, Y=6: '╟'} \\
\texttt{X=2, Y=6: '─'} \\
\texttt{X=3, Y=6: '─'} \\
\texttt{X=4, Y=6: '─'} \\
\texttt{X=5, Y=6: '┼'} \\
\texttt{X=6, Y=6: '─'} \\
\texttt{X=7, Y=6: '┘'} \\
\\ 
\texttt{Да, (1,6) занято.} \\
\texttt{Значит, цель A: (1,6) недостижима.} \\
\texttt{Нам нужно договориться о новой цели для A.} \\
\texttt{Или, возможно, я должен был двигаться по другому пути.} \\
\\ 
\texttt{Давай предложим новую стратегию.} \\
\texttt{Поскольку (1,6) занято, давай перенесем цель A в (1,5)? Но там пусто.} \\
\texttt{Или (2,6)? Занято.} \\
\texttt{(3,6)? Занято.} \\
\texttt{(4,6)? Занято.} \\
\texttt{(5,6)? Занято.} \\
\texttt{(6,6)? Занято.} \\
\texttt{(7,6)? Занято (там A сейчас).} \\
\\ 
\texttt{Давай посмотрим на строку 5:} \\
\texttt{║◌◌◌│◌├─╢ 5} \\
\texttt{X=1: ◌} \\
\texttt{X=2: ◌} \\
\texttt{X=3: ◌} \\
\texttt{X=4: │} \\
\texttt{X=5: ◌} \\
\texttt{X=6: ├} \\
\texttt{X=7: ─} \\
\\ 
\texttt{(1,5) пусто.} \\
\texttt{Давай перенесем цель A в (1,5).} \\
\texttt{Но тогда D, который должен быть в (1,5), должен куда{-}то деться.} \\
\texttt{D цель (1,5).} \\
\texttt{Если A в (1,5), то D должен быть в другом месте.} \\
\\ 
\texttt{Давай договоримся о новой конфигурации.} \\
\texttt{P: (1,1)} \\
\texttt{O: (1,2)} \\
\texttt{T: (1,3)} \\
\texttt{S: (1,4)} \\
\texttt{D: (1,5)} \\
\texttt{A: (2,5) ?} \\
\texttt{M: (1,7)} \\
\\ 
\texttt{Но (2,5) пусто?} \\
\texttt{Строка 5, столбец 2: '◌'. Да, пусто.} \\
\\ 
\texttt{Давай предложим это.} \\
\\ 
\texttt{СКАЗАТЬ: Столбец 1 в строке 6 заблокирован линией, поэтому цель A: (1,6) недостижима. Давай изменим цели для A и D.} \\
\texttt{Новые цели:} \\
\texttt{P: (1,1)} \\
\texttt{O: (1,2)} \\
\texttt{T: (1,3)} \\
\texttt{S: (1,4)} \\
\texttt{D: (1,5)} \\
\texttt{A: (2,5)} \\
\texttt{M: (1,7)} \\
\\ 
\texttt{Если ты согласен, давай двигаться.} \\
\texttt{Сначала переместим M из (1,3) в (1,7)? Нет, там D.} \\
\texttt{Сначала переместим D из (1,7) в (1,5).} \\
\texttt{Путь: (1,7) {-}> (1,6) занято.} \\
\texttt{Нуж} \\
            }
        }
    }
     \\ \\

    \theutterance \stepcounter{utterance}  
    & & & \multicolumn{2}{p{0.3\linewidth}}{
        \cellcolor[rgb]{0.9,0.9,0.9}{
            \makecell[{{p{\linewidth}}}]{
                \texttt{\tiny{[GM$|$GM]}}
                \texttt{СКАЗАТЬ: Мы накопили много штрафов, пытаясь двигаться по прямой линии из{-}за препятствий. Давай изменим стратегию. Вместо того чтобы двигаться строго по столбцу 1, давай двигаться к цели по кратчайшему пути, избегая препятствий.} \\
\\ 
\texttt{Текущие цели:} \\
\texttt{P: (1, 1)} \\
\texttt{O: (1, 2) {-} уже на месте} \\
\texttt{T: (1, 3)} \\
\texttt{S: (1, 4) {-} уже на месте} \\
\texttt{D: (1, 5)} \\
\texttt{A: (1, 6) {-} сейчас в (7,6), нужно вернуть в (1,6)} \\
\texttt{M: (1, 7) {-} сейчас в (1,3), нужно переместить в (1,7)} \\
\\ 
\texttt{Проблема в том, что столбец 1 занят рамкой и линиями внизу. Нам нужно обходить препятствия.} \\
\texttt{Давай договоримся о следующем плане действий для оставшихся объектов:} \\
\texttt{1. M сейчас в (1,3). Цель (1,7). Путь вниз по столбцу 1 заблокирован линиями в (5,1) и (6,1). Нужно обходить.} \\
\texttt{2. D сейчас в (1,7). Цель (1,5). Путь вверх по столбцу 1 заблокирован.} \\
\texttt{3. T сейчас в (1,7). Цель (1,3).} \\
\texttt{4. P сейчас в (1,7). Цель (1,1).} \\
\texttt{5. A сейчас в (7,6). Цель (1,6).} \\
\\ 
\texttt{Давай попробуем переместить объекты, которые могут двигаться свободно, чтобы освободить место или создать путь.} \\
\texttt{Сначала давай переместим M из (1,3) в (2,3), чтобы освободить (1,3) для T, если нужно, или просто сдвинуть.} \\
\texttt{Но лучше давай сначала переместим P, T, D, которые находятся в (1,7).} \\
\texttt{P цель (1,1). T цель (1,3). D цель (1,5).} \\
\\ 
\texttt{Давай начнем с перемещения P из (1,7) в (1,6)? Нет, там A.} \\
\texttt{Давай переместим A из (7,6) в (1,6). Путь: (7,6) {-}> (6,6) занято, (5,6) занято...} \\
\texttt{Сетка:} \\
\texttt{(1,1) пусто} \\
\texttt{(1,2) O} \\
\texttt{(1,3) M} \\
\texttt{(1,4) S} \\
\texttt{(1,5) пусто? Нет, в (5,1) есть линия.} \\
\texttt{(1,6) пусто? Нет, в (6,1) есть линия.} \\
\texttt{(1,7) пусто? Нет, там D, T, P.} \\
\\ 
\texttt{Посмотрим на столбец 1:} \\
\texttt{(1,1) пусто} \\
\texttt{(1,2) O} \\
\texttt{(1,3) M} \\
\texttt{(1,4) S} \\
\texttt{(1,5) пусто (в (5,1) линия, но (1,5) это строка 5, столбец 1. В (5,1) есть '|'. Значит (1,5) занято линией? Нет, координаты (X,Y). X=1, Y=5. В сетке:} \\
\texttt{║◌◌◌│◌├─╢ 5} \\
\texttt{Столбец 1, строка 5: '◌'. Пусто.} \\
\texttt{Строка 6:} \\
\texttt{╟───┼─┘◌║ 6} \\
\texttt{Столбец 1, строка 6: '╟'. Занято.} \\
\texttt{Строка 7:} \\
\texttt{║D◌T│P◌◌║ 7} \\
\texttt{Столбец 1, строка 7: 'D'. Занято.} \\
\\ 
\texttt{Итак, в столбце 1:} \\
\texttt{(1,1) пусто} \\
\texttt{(1,2) O} \\
\texttt{(1,3) M} \\
\texttt{(1,4) S} \\
\texttt{(1,5) пусто} \\
\texttt{(1,6) пусто (в (6,1) '╟', но (1,6) это столбец 1, строка 6. В строке 6 столбец 1 это '╟'. Занято.)} \\
\texttt{(1,7) D} \\
\\ 
\texttt{Давай проверим координаты еще раз.} \\
\texttt{X {-} столбец, Y {-} строка.} \\
\texttt{(1,1) {-} пусто} \\
\texttt{(1,2) {-} O} \\
\texttt{(1,3) {-} M} \\
\texttt{(1,4) {-} S} \\
\texttt{(1,5) {-} пусто (в строке 5, столбец 1 это '◌')} \\
\texttt{(1,6) {-} занято ('╟' в строке 6, столбец 1)} \\
\texttt{(1,7) {-} D} \\
\\ 
\texttt{Цели:} \\
\texttt{P: (1,1)} \\
\texttt{O: (1,2) {-} на месте} \\
\texttt{T: (1,3)} \\
\texttt{S: (1,4) {-} на месте} \\
\texttt{D: (1,5)} \\
\texttt{A: (1,6)} \\
\texttt{M: (1,7)} \\
\\ 
\texttt{Проблема: (1,6) занято линией. Мы не можем поместить A в (1,6).} \\
\texttt{Значит, наша целевая конфигурация невозможна в точных координатах, если (1,6) заблокировано.} \\
\texttt{Но в начале игры мы договорились о (1,6) для A.} \\
\texttt{Возможно, я неправильно понял структуру сетки.} \\
\texttt{Давай посмотрим на строку 6:} \\
\texttt{╟───┼─┘◌║ 6} \\
\texttt{Символы:} \\
\texttt{1: ╟} \\
\texttt{2: ─} \\
\texttt{3: ─} \\
\texttt{4: ─} \\
\texttt{5: ┼} \\
\texttt{6: ─} \\
\texttt{7: ┘} \\
\texttt{8: ◌} \\
\texttt{9: ║} \\
\\ 
\texttt{Координаты X от 1 до 7.} \\
\texttt{X=1, Y=6: '╟'} \\
\texttt{X=2, Y=6: '─'} \\
\texttt{X=3, Y=6: '─'} \\
\texttt{X=4, Y=6: '─'} \\
\texttt{X=5, Y=6: '┼'} \\
\texttt{X=6, Y=6: '─'} \\
\texttt{X=7, Y=6: '┘'} \\
\\ 
\texttt{Да, (1,6) занято.} \\
\texttt{Значит, цель A: (1,6) недостижима.} \\
\texttt{Нам нужно договориться о новой цели для A.} \\
\texttt{Или, возможно, я должен был двигаться по другому пути.} \\
\\ 
\texttt{Давай предложим новую стратегию.} \\
\texttt{Поскольку (1,6) занято, давай перенесем цель A в (1,5)? Но там пусто.} \\
\texttt{Или (2,6)? Занято.} \\
\texttt{(3,6)? Занято.} \\
\texttt{(4,6)? Занято.} \\
\texttt{(5,6)? Занято.} \\
\texttt{(6,6)? Занято.} \\
\texttt{(7,6)? Занято (там A сейчас).} \\
\\ 
\texttt{Давай посмотрим на строку 5:} \\
\texttt{║◌◌◌│◌├─╢ 5} \\
\texttt{X=1: ◌} \\
\texttt{X=2: ◌} \\
\texttt{X=3: ◌} \\
\texttt{X=4: │} \\
\texttt{X=5: ◌} \\
\texttt{X=6: ├} \\
\texttt{X=7: ─} \\
\\ 
\texttt{(1,5) пусто.} \\
\texttt{Давай перенесем цель A в (1,5).} \\
\texttt{Но тогда D, который должен быть в (1,5), должен куда{-}то деться.} \\
\texttt{D цель (1,5).} \\
\texttt{Если A в (1,5), то D должен быть в другом месте.} \\
\\ 
\texttt{Давай договоримся о новой конфигурации.} \\
\texttt{P: (1,1)} \\
\texttt{O: (1,2)} \\
\texttt{T: (1,3)} \\
\texttt{S: (1,4)} \\
\texttt{D: (1,5)} \\
\texttt{A: (2,5) ?} \\
\texttt{M: (1,7)} \\
\\ 
\texttt{Но (2,5) пусто?} \\
\texttt{Строка 5, столбец 2: '◌'. Да, пусто.} \\
\\ 
\texttt{Давай предложим это.} \\
\\ 
\texttt{СКАЗАТЬ: Столбец 1 в строке 6 заблокирован линией, поэтому цель A: (1,6) недостижима. Давай изменим цели для A и D.} \\
\texttt{Новые цели:} \\
\texttt{P: (1,1)} \\
\texttt{O: (1,2)} \\
\texttt{T: (1,3)} \\
\texttt{S: (1,4)} \\
\texttt{D: (1,5)} \\
\texttt{A: (2,5)} \\
\texttt{M: (1,7)} \\
\\ 
\texttt{Если ты согласен, давай двигаться.} \\
\texttt{Сначала переместим M из (1,3) в (1,7)? Нет, там D.} \\
\texttt{Сначала переместим D из (1,7) в (1,5).} \\
\texttt{Путь: (1,7) {-}> (1,6) занято.} \\
\texttt{Нуж} \\
            }
        }
    }
    & & \\ \\

    \theutterance \stepcounter{utterance}  
    & & & \multicolumn{2}{p{0.3\linewidth}}{
        \cellcolor[rgb]{0.9,0.9,0.9}{
            \makecell[{{p{\linewidth}}}]{
                \texttt{\tiny{[GM$|$GM]}}
                \texttt{Invalid format: head and tail are not empty} \\
\texttt{head: 'СКАЗАТЬ: Мы накопили много штрафов, пытаясь двигаться по прямой линии из{-}за препятствий. Давай изменим стратегию. Вместо того чтобы двигаться строго по столбцу 1, давай двигаться к цели по кратчайшему пути, избегая препятствий.} \\
\\ 
\texttt{Текущие цели:} \\
\texttt{P: (1, 1)} \\
\texttt{O: (1, 2) {-} уже на месте} \\
\texttt{T: (1, 3)} \\
\texttt{S: (1, 4) {-} уже на месте} \\
\texttt{D: (1, 5)} \\
\texttt{A: (1, 6) {-} сейчас в (7,6), нужно вернуть в (1,6)} \\
\texttt{M: (1, 7) {-} сейчас в (1,3), нужно переместить в (1,7)} \\
\\ 
\texttt{Проблема в том, что столбец 1 занят рамкой и линиями внизу. Нам нужно обходить препятствия.} \\
\texttt{Давай договоримся о следующем плане действий для оставшихся объектов:} \\
\texttt{1. M сейчас в (1,3). Цель (1,7). Путь вниз по столбцу 1 заблокирован линиями в (5,1) и (6,1). Нужно обходить.} \\
\texttt{2. D сейчас в (1,7). Цель (1,5). Путь вверх по столбцу 1 заблокирован.} \\
\texttt{3. T сейчас в (1,7). Цель (1,3).} \\
\texttt{4. P сейчас в (1,7). Цель (1,1).} \\
\texttt{5. A сейчас в (7,6). Цель (1,6).} \\
\\ 
\texttt{Давай попробуем переместить объекты, которые могут двигаться свободно, чтобы освободить место или создать путь.} \\
\texttt{Сначала давай переместим M из (1,3) в (2,3), чтобы освободить (1,3) для T, если нужно, или просто сдвинуть.} \\
\texttt{Но лучше давай сначала переместим P, T, D, которые находятся в (1,7).} \\
\texttt{P цель (1,1). T цель (1,3). D цель (1,5).} \\
\\ 
\texttt{Давай начнем с перемещения P из (1,7) в (1,6)? Нет, там A.} \\
\texttt{Давай переместим A из (7,6) в (1,6). Путь: (7,6) {-}> (6,6) занято, (5,6) занято...} \\
\texttt{Сетка:} \\
\texttt{(1,1) пусто} \\
\texttt{(1,2) O} \\
\texttt{(1,3) M} \\
\texttt{(1,4) S} \\
\texttt{(1,5) пусто? Нет, в (5,1) есть линия.} \\
\texttt{(1,6) пусто? Нет, в (6,1) есть линия.} \\
\texttt{(1,7) пусто? Нет, там D, T, P.} \\
\\ 
\texttt{Посмотрим на столбец 1:} \\
\texttt{(1,1) пусто} \\
\texttt{(1,2) O} \\
\texttt{(1,3) M} \\
\texttt{(1,4) S} \\
\texttt{(1,5) пусто (в (5,1) линия, но (1,5) это строка 5, столбец 1. В (5,1) есть '|'. Значит (1,5) занято линией? Нет, координаты (X,Y). X=1, Y=5. В сетке:} \\
\texttt{║◌◌◌│◌├─╢ 5} \\
\texttt{Столбец 1, строка 5: '◌'. Пусто.} \\
\texttt{Строка 6:} \\
\texttt{╟───┼─┘◌║ 6} \\
\texttt{Столбец 1, строка 6: '╟'. Занято.} \\
\texttt{Строка 7:} \\
\texttt{║D◌T│P◌◌║ 7} \\
\texttt{Столбец 1, строка 7: 'D'. Занято.} \\
\\ 
\texttt{Итак, в столбце 1:} \\
\texttt{(1,1) пусто} \\
\texttt{(1,2) O} \\
\texttt{(1,3) M} \\
\texttt{(1,4) S} \\
\texttt{(1,5) пусто} \\
\texttt{(1,6) пусто (в (6,1) '╟', но (1,6) это столбец 1, строка 6. В строке 6 столбец 1 это '╟'. Занято.)} \\
\texttt{(1,7) D} \\
\\ 
\texttt{Давай проверим координаты еще раз.} \\
\texttt{X {-} столбец, Y {-} строка.} \\
\texttt{(1,1) {-} пусто} \\
\texttt{(1,2) {-} O} \\
\texttt{(1,3) {-} M} \\
\texttt{(1,4) {-} S} \\
\texttt{(1,5) {-} пусто (в строке 5, столбец 1 это '◌')} \\
\texttt{(1,6) {-} занято ('╟' в строке 6, столбец 1)} \\
\texttt{(1,7) {-} D} \\
\\ 
\texttt{Цели:} \\
\texttt{P: (1,1)} \\
\texttt{O: (1,2) {-} на месте} \\
\texttt{T: (1,3)} \\
\texttt{S: (1,4) {-} на месте} \\
\texttt{D: (1,5)} \\
\texttt{A: (1,6)} \\
\texttt{M: (1,7)} \\
\\ 
\texttt{Проблема: (1,6) занято линией. Мы не можем поместить A в (1,6).} \\
\texttt{Значит, наша целевая конфигурация невозможна в точных координатах, если (1,6) заблокировано.} \\
\texttt{Но в начале игры мы договорились о (1,6) для A.} \\
\texttt{Возможно, я неправильно понял структуру сетки.} \\
\texttt{Давай посмотрим на строку 6:} \\
\texttt{╟───┼─┘◌║ 6} \\
\texttt{Символы:} \\
\texttt{1: ╟} \\
\texttt{2: ─} \\
\texttt{3: ─} \\
\texttt{4: ─} \\
\texttt{5: ┼} \\
\texttt{6: ─} \\
\texttt{7: ┘} \\
\texttt{8: ◌} \\
\texttt{9: ║} \\
\\ 
\texttt{Координаты X от 1 до 7.} \\
\texttt{X=1, Y=6: '╟'} \\
\texttt{X=2, Y=6: '─'} \\
\texttt{X=3, Y=6: '─'} \\
\texttt{X=4, Y=6: '─'} \\
\texttt{X=5, Y=6: '┼'} \\
\texttt{X=6, Y=6: '─'} \\
\texttt{X=7, Y=6: '┘'} \\
\\ 
\texttt{Да, (1,6) занято.} \\
\texttt{Значит, цель A: (1,6) недостижима.} \\
\texttt{Нам нужно договориться о новой цели для A.} \\
\texttt{Или, возможно, я должен был двигаться по другому пути.} \\
\\ 
\texttt{Давай предложим новую стратегию.} \\
\texttt{Поскольку (1,6) занято, давай перенесем цель A в (1,5)? Но там пусто.} \\
\texttt{Или (2,6)? Занято.} \\
\texttt{(3,6)? Занято.} \\
\texttt{(4,6)? Занято.} \\
\texttt{(5,6)? Занято.} \\
\texttt{(6,6)? Занято.} \\
\texttt{(7,6)? Занято (там A сейчас).} \\
\\ 
\texttt{Давай посмотрим на строку 5:} \\
\texttt{║◌◌◌│◌├─╢ 5} \\
\texttt{X=1: ◌} \\
\texttt{X=2: ◌} \\
\texttt{X=3: ◌} \\
\texttt{X=4: │} \\
\texttt{X=5: ◌} \\
\texttt{X=6: ├} \\
\texttt{X=7: ─} \\
\\ 
\texttt{(1,5) пусто.} \\
\texttt{Давай перенесем цель A в (1,5).} \\
\texttt{Но тогда D, который должен быть в (1,5), должен куда{-}то деться.} \\
\texttt{D цель (1,5).} \\
\texttt{Если A в (1,5), то D должен быть в другом месте.} \\
\\ 
\texttt{Давай договоримся о новой конфигурации.} \\
\texttt{P: (1,1)} \\
\texttt{O: (1,2)} \\
\texttt{T: (1,3)} \\
\texttt{S: (1,4)} \\
\texttt{D: (1,5)} \\
\texttt{A: (2,5) ?} \\
\texttt{M: (1,7)} \\
\\ 
\texttt{Но (2,5) пусто?} \\
\texttt{Строка 5, столбец 2: '◌'. Да, пусто.} \\
\\ 
\texttt{Давай предложим это.} \\
\\ 
\texttt{'} \\
\texttt{tail: '} \\
\texttt{Новые цели:} \\
\texttt{P: (1,1)} \\
\texttt{O: (1,2)} \\
\texttt{T: (1,3)} \\
\texttt{S: (1,4)} \\
\texttt{D: (1,5)} \\
\texttt{A: (2,5)} \\
\texttt{M: (1,7)} \\
\\ 
\texttt{Если ты согласен, давай двигаться.} \\
\texttt{Сначала переместим M из (1,3) в (1,7)? Нет, там D.} \\
\texttt{Сначала переместим D из (1,7) в (1,5).} \\
\texttt{Путь: (1,7) {-}> (1,6) занято.} \\
\texttt{Нуж'} \\
            }
        }
    }
    & & \\ \\

    \theutterance \stepcounter{utterance}  
    & & \multicolumn{4}{p{0.6\linewidth}}{
        \cellcolor[rgb]{0.9,0.9,0.9}{
            \makecell[{{p{\linewidth}}}]{
                \texttt{\tiny{[GM$\rangle$P2]}}
                \texttt{Штраф: Ваше сообщение не должно содержать ничего ни до, ни после команды!} \\
\\ 
\texttt{Убедитесь, что ваш ответ содержит только `СКАЗАТЬ: <СООБЩЕНИЕ>` или `ДВИГАТЬСЯ: <ОБЪЕКТ>, (<X>, <Y>)` и ничего больше!} \\
\\ 
\texttt{Вы вместе накопили 10 из 16 штрафов.} \\
\texttt{Пожалуйста, попробуйте ещё раз!} \\
            }
        }
    }
    & \\ \\

    \theutterance \stepcounter{utterance}  
    & & & \multicolumn{4}{p{0.6\linewidth}}{
        \cellcolor[rgb]{0.9,0.9,0.9}{
            \makecell[{{p{\linewidth}}}]{
                \texttt{\tiny{[GM$\langle$P2]}}
                \texttt{СКАЗАТЬ: Столбец 1 в строке 6 заблокирован, поэтому цель A:(1,6) недостижима. Предлагаю новые цели: P:(1,1), O:(1,2), T:(1,3), S:(1,4), D:(1,5), A:(2,5), M:(1,7). Согласен?} \\
            }
        }
    }
     \\ \\

    \theutterance \stepcounter{utterance}  
    & & & \multicolumn{2}{p{0.3\linewidth}}{
        \cellcolor[rgb]{0.9,0.9,0.9}{
            \makecell[{{p{\linewidth}}}]{
                \texttt{\tiny{[GM$|$GM]}}
                \texttt{СКАЗАТЬ: Столбец 1 в строке 6 заблокирован, поэтому цель A:(1,6) недостижима. Предлагаю новые цели: P:(1,1), O:(1,2), T:(1,3), S:(1,4), D:(1,5), A:(2,5), M:(1,7). Согласен?} \\
            }
        }
    }
    & & \\ \\

    \theutterance \stepcounter{utterance}  
    & & \multicolumn{4}{p{0.6\linewidth}}{
        \cellcolor[rgb]{0.9,0.9,0.9}{
            \makecell[{{p{\linewidth}}}]{
                \texttt{\tiny{[P1$\langle$GM]}}
                \texttt{Успешно перемещён объект "D" в (1,5). Ваша обновлённая сетка выглядит так:} \\
\texttt{```} \\
\texttt{ 1234567} \\
\texttt{╔═══════╗} \\
\texttt{║P◌◌◌◌◌◌║ 1} \\
\texttt{║O◌◌◌◌◌◌║ 2} \\
\texttt{║T◌A◌◌┌─╢ 3} \\
\texttt{║S◌◌┌─┤◌║ 4} \\
\texttt{║D◌◌│◌├─╢ 5} \\
\texttt{╟───┼─┘◌║ 6} \\
\texttt{║◌M◌│◌◌◌║ 7} \\
\texttt{╚═══╧═══╝} \\
\texttt{```} \\
\texttt{Сейчас вы играете раунд 6 из максимум 28.} \\
\texttt{Вы вместе накопили 10 из 16 штрафов.} \\
\texttt{Другой игрок отправил это сообщение:} \\
\\ 
\\ 
\texttt{"Столбец 1 в строке 6 заблокирован, поэтому цель A:(1,6) недостижима. Предлагаю новые цели: P:(1,1), O:(1,2), T:(1,3), S:(1,4), D:(1,5), A:(2,5), M:(1,7). Согласен?"} \\
\\ 
\\ 
\texttt{Какая твоя следующая команда?} \\
            }
        }
    }
    & \\ \\

    \theutterance \stepcounter{utterance}  
    & \multicolumn{4}{p{0.6\linewidth}}{
        \cellcolor[rgb]{0.9,0.9,0.9}{
            \makecell[{{p{\linewidth}}}]{
                \texttt{\tiny{[P1$\rangle$GM]}}
                \texttt{СКАЗАТЬ: Согласен с новым планом. Перемещаем A в (2,5).} \\
            }
        }
    }
    & & \\ \\

    \theutterance \stepcounter{utterance}  
    & & & \multicolumn{2}{p{0.3\linewidth}}{
        \cellcolor[rgb]{0.9,0.9,0.9}{
            \makecell[{{p{\linewidth}}}]{
                \texttt{\tiny{[GM$|$GM]}}
                \texttt{СКАЗАТЬ: Согласен с новым планом. Перемещаем A в (2,5).} \\
            }
        }
    }
    & & \\ \\

    \theutterance \stepcounter{utterance}  
    & & \multicolumn{4}{p{0.6\linewidth}}{
        \cellcolor[rgb]{0.9,0.9,0.9}{
            \makecell[{{p{\linewidth}}}]{
                \texttt{\tiny{[GM$\rangle$P2]}}
                \texttt{Ваше сообщение было передано другому игроку.} \\
\\ 
\\ 
\texttt{Сейчас вы играете раунд 7 из максимум 28.} \\
\texttt{Вы вместе накопили 10 из 16 штрафов.} \\
\texttt{Другой игрок отправил это сообщение:} \\
\\ 
\\ 
\texttt{"Согласен с новым планом. Перемещаем A в (2,5)."} \\
\\ 
\\ 
\texttt{Какая твоя следующая команда?} \\
            }
        }
    }
    & \\ \\

    \theutterance \stepcounter{utterance}  
    & & & \multicolumn{4}{p{0.6\linewidth}}{
        \cellcolor[rgb]{0.9,0.9,0.9}{
            \makecell[{{p{\linewidth}}}]{
                \texttt{\tiny{[GM$\langle$P2]}}
                \texttt{ДВИГАТЬСЯ: A, (2, 5)} \\
            }
        }
    }
     \\ \\

    \theutterance \stepcounter{utterance}  
    & & & \multicolumn{2}{p{0.3\linewidth}}{
        \cellcolor[rgb]{0.9,0.9,0.9}{
            \makecell[{{p{\linewidth}}}]{
                \texttt{\tiny{[GM$|$GM]}}
                \texttt{ДВИГАТЬСЯ: A, (2, 5)} \\
            }
        }
    }
    & & \\ \\

    \theutterance \stepcounter{utterance}  
    & & & \multicolumn{2}{p{0.3\linewidth}}{
        \cellcolor[rgb]{0.9,0.9,0.9}{
            \makecell[{{p{\linewidth}}}]{
                \texttt{\tiny{[GM$|$GM]}}
                \texttt{A moved to (2, 5)} \\
            }
        }
    }
    & & \\ \\

    \theutterance \stepcounter{utterance}  
    & & \multicolumn{4}{p{0.6\linewidth}}{
        \cellcolor[rgb]{0.9,0.9,0.9}{
            \makecell[{{p{\linewidth}}}]{
                \texttt{\tiny{[P1$\langle$GM]}}
                \texttt{Ваше сообщение было передано другому игроку.} \\
\\ 
\\ 
\texttt{Сейчас вы играете раунд 7 из максимум 28.} \\
\texttt{Вы вместе накопили 10 из 16 штрафов.} \\
\texttt{Другой игрок переместил объект на своей сетке.} \\
\\ 
\\ 
\texttt{Какая твоя следующая команда?} \\
            }
        }
    }
    & \\ \\

    \theutterance \stepcounter{utterance}  
    & \multicolumn{4}{p{0.6\linewidth}}{
        \cellcolor[rgb]{0.9,0.9,0.9}{
            \makecell[{{p{\linewidth}}}]{
                \texttt{\tiny{[P1$\rangle$GM]}}
                \texttt{ДВИГАТЬСЯ: A, (2, 3)} \\
            }
        }
    }
    & & \\ \\

    \theutterance \stepcounter{utterance}  
    & & & \multicolumn{2}{p{0.3\linewidth}}{
        \cellcolor[rgb]{0.9,0.9,0.9}{
            \makecell[{{p{\linewidth}}}]{
                \texttt{\tiny{[GM$|$GM]}}
                \texttt{ДВИГАТЬСЯ: A, (2, 3)} \\
            }
        }
    }
    & & \\ \\

    \theutterance \stepcounter{utterance}  
    & & & \multicolumn{2}{p{0.3\linewidth}}{
        \cellcolor[rgb]{0.9,0.9,0.9}{
            \makecell[{{p{\linewidth}}}]{
                \texttt{\tiny{[GM$|$GM]}}
                \texttt{A moved to (2, 3)} \\
            }
        }
    }
    & & \\ \\

    \theutterance \stepcounter{utterance}  
    & & \multicolumn{4}{p{0.6\linewidth}}{
        \cellcolor[rgb]{0.9,0.9,0.9}{
            \makecell[{{p{\linewidth}}}]{
                \texttt{\tiny{[GM$\rangle$P2]}}
                \texttt{Успешно перемещён объект "A" в (2,5). Ваша обновлённая сетка выглядит так:} \\
\texttt{```} \\
\texttt{ 1234567} \\
\texttt{╔═══════╗} \\
\texttt{║◌◌◌◌◌◌◌║ 1} \\
\texttt{║O◌◌◌◌◌◌║ 2} \\
\texttt{║M◌◌◌◌┌─╢ 3} \\
\texttt{║S◌◌┌─┤◌║ 4} \\
\texttt{║◌A◌│◌├─╢ 5} \\
\texttt{╟───┼─┘◌║ 6} \\
\texttt{║D◌T│P◌◌║ 7} \\
\texttt{╚═══╧═══╝} \\
\texttt{```} \\
\texttt{Сейчас вы играете раунд 8 из максимум 28.} \\
\texttt{Вы вместе накопили 10 из 16 штрафов.} \\
\texttt{Другой игрок переместил объект на своей сетке.} \\
\\ 
\\ 
\texttt{Какая твоя следующая команда?} \\
            }
        }
    }
    & \\ \\

    \theutterance \stepcounter{utterance}  
    & & & \multicolumn{4}{p{0.6\linewidth}}{
        \cellcolor[rgb]{0.9,0.9,0.9}{
            \makecell[{{p{\linewidth}}}]{
                \texttt{\tiny{[GM$\langle$P2]}}
                \texttt{ДВИГАТЬСЯ: D, (1, 5)} \\
            }
        }
    }
     \\ \\

    \theutterance \stepcounter{utterance}  
    & & & \multicolumn{2}{p{0.3\linewidth}}{
        \cellcolor[rgb]{0.9,0.9,0.9}{
            \makecell[{{p{\linewidth}}}]{
                \texttt{\tiny{[GM$|$GM]}}
                \texttt{ДВИГАТЬСЯ: D, (1, 5)} \\
            }
        }
    }
    & & \\ \\

    \theutterance \stepcounter{utterance}  
    & & & \multicolumn{2}{p{0.3\linewidth}}{
        \cellcolor[rgb]{0.9,0.9,0.9}{
            \makecell[{{p{\linewidth}}}]{
                \texttt{\tiny{[GM$|$GM]}}
                \texttt{D moved to (1, 5)} \\
            }
        }
    }
    & & \\ \\

    \theutterance \stepcounter{utterance}  
    & & \multicolumn{4}{p{0.6\linewidth}}{
        \cellcolor[rgb]{0.9,0.9,0.9}{
            \makecell[{{p{\linewidth}}}]{
                \texttt{\tiny{[P1$\langle$GM]}}
                \texttt{Успешно перемещён объект "A" в (2,3). Ваша обновлённая сетка выглядит так:} \\
\texttt{```} \\
\texttt{ 1234567} \\
\texttt{╔═══════╗} \\
\texttt{║P◌◌◌◌◌◌║ 1} \\
\texttt{║O◌◌◌◌◌◌║ 2} \\
\texttt{║TA◌◌◌┌─╢ 3} \\
\texttt{║S◌◌┌─┤◌║ 4} \\
\texttt{║D◌◌│◌├─╢ 5} \\
\texttt{╟───┼─┘◌║ 6} \\
\texttt{║◌M◌│◌◌◌║ 7} \\
\texttt{╚═══╧═══╝} \\
\texttt{```} \\
\texttt{Сейчас вы играете раунд 8 из максимум 28.} \\
\texttt{Вы вместе накопили 10 из 16 штрафов.} \\
\texttt{Другой игрок переместил объект на своей сетке.} \\
\\ 
\\ 
\texttt{Какая твоя следующая команда?} \\
            }
        }
    }
    & \\ \\

    \theutterance \stepcounter{utterance}  
    & \multicolumn{4}{p{0.6\linewidth}}{
        \cellcolor[rgb]{0.9,0.9,0.9}{
            \makecell[{{p{\linewidth}}}]{
                \texttt{\tiny{[P1$\rangle$GM]}}
                \texttt{ДВИГАТЬСЯ: A, (2, 5)} \\
            }
        }
    }
    & & \\ \\

    \theutterance \stepcounter{utterance}  
    & & & \multicolumn{2}{p{0.3\linewidth}}{
        \cellcolor[rgb]{0.9,0.9,0.9}{
            \makecell[{{p{\linewidth}}}]{
                \texttt{\tiny{[GM$|$GM]}}
                \texttt{ДВИГАТЬСЯ: A, (2, 5)} \\
            }
        }
    }
    & & \\ \\

    \theutterance \stepcounter{utterance}  
    & & & \multicolumn{2}{p{0.3\linewidth}}{
        \cellcolor[rgb]{0.9,0.9,0.9}{
            \makecell[{{p{\linewidth}}}]{
                \texttt{\tiny{[GM$|$GM]}}
                \texttt{A moved to (2, 5)} \\
            }
        }
    }
    & & \\ \\

    \theutterance \stepcounter{utterance}  
    & & \multicolumn{4}{p{0.6\linewidth}}{
        \cellcolor[rgb]{0.9,0.9,0.9}{
            \makecell[{{p{\linewidth}}}]{
                \texttt{\tiny{[GM$\rangle$P2]}}
                \texttt{Успешно перемещён объект "D" в (1,5). Ваша обновлённая сетка выглядит так:} \\
\texttt{```} \\
\texttt{ 1234567} \\
\texttt{╔═══════╗} \\
\texttt{║◌◌◌◌◌◌◌║ 1} \\
\texttt{║O◌◌◌◌◌◌║ 2} \\
\texttt{║M◌◌◌◌┌─╢ 3} \\
\texttt{║S◌◌┌─┤◌║ 4} \\
\texttt{║DA◌│◌├─╢ 5} \\
\texttt{╟───┼─┘◌║ 6} \\
\texttt{║◌◌T│P◌◌║ 7} \\
\texttt{╚═══╧═══╝} \\
\texttt{```} \\
\texttt{Сейчас вы играете раунд 9 из максимум 28.} \\
\texttt{Вы вместе накопили 10 из 16 штрафов.} \\
\texttt{Другой игрок переместил объект на своей сетке.} \\
\\ 
\\ 
\texttt{Какая твоя следующая команда?} \\
            }
        }
    }
    & \\ \\

    \theutterance \stepcounter{utterance}  
    & & & \multicolumn{4}{p{0.6\linewidth}}{
        \cellcolor[rgb]{0.9,0.9,0.9}{
            \makecell[{{p{\linewidth}}}]{
                \texttt{\tiny{[GM$\langle$P2]}}
                \texttt{ДВИГАТЬСЯ: M, (1, 7)} \\
            }
        }
    }
     \\ \\

    \theutterance \stepcounter{utterance}  
    & & & \multicolumn{2}{p{0.3\linewidth}}{
        \cellcolor[rgb]{0.9,0.9,0.9}{
            \makecell[{{p{\linewidth}}}]{
                \texttt{\tiny{[GM$|$GM]}}
                \texttt{ДВИГАТЬСЯ: M, (1, 7)} \\
            }
        }
    }
    & & \\ \\

    \theutterance \stepcounter{utterance}  
    & & & \multicolumn{2}{p{0.3\linewidth}}{
        \cellcolor[rgb]{0.9,0.9,0.9}{
            \makecell[{{p{\linewidth}}}]{
                \texttt{\tiny{[GM$|$GM]}}
                \texttt{M moved to (1, 7)} \\
            }
        }
    }
    & & \\ \\

    \theutterance \stepcounter{utterance}  
    & & \multicolumn{4}{p{0.6\linewidth}}{
        \cellcolor[rgb]{0.9,0.9,0.9}{
            \makecell[{{p{\linewidth}}}]{
                \texttt{\tiny{[P1$\langle$GM]}}
                \texttt{Успешно перемещён объект "A" в (2,5). Ваша обновлённая сетка выглядит так:} \\
\texttt{```} \\
\texttt{ 1234567} \\
\texttt{╔═══════╗} \\
\texttt{║P◌◌◌◌◌◌║ 1} \\
\texttt{║O◌◌◌◌◌◌║ 2} \\
\texttt{║T◌◌◌◌┌─╢ 3} \\
\texttt{║S◌◌┌─┤◌║ 4} \\
\texttt{║DA◌│◌├─╢ 5} \\
\texttt{╟───┼─┘◌║ 6} \\
\texttt{║◌M◌│◌◌◌║ 7} \\
\texttt{╚═══╧═══╝} \\
\texttt{```} \\
\texttt{Сейчас вы играете раунд 9 из максимум 28.} \\
\texttt{Вы вместе накопили 10 из 16 штрафов.} \\
\texttt{Другой игрок переместил объект на своей сетке.} \\
\\ 
\\ 
\texttt{Какая твоя следующая команда?} \\
            }
        }
    }
    & \\ \\

    \theutterance \stepcounter{utterance}  
    & \multicolumn{4}{p{0.6\linewidth}}{
        \cellcolor[rgb]{0.9,0.9,0.9}{
            \makecell[{{p{\linewidth}}}]{
                \texttt{\tiny{[P1$\rangle$GM]}}
                \texttt{ДВИГАТЬСЯ: M, (1, 7)} \\
            }
        }
    }
    & & \\ \\

    \theutterance \stepcounter{utterance}  
    & & & \multicolumn{2}{p{0.3\linewidth}}{
        \cellcolor[rgb]{0.9,0.9,0.9}{
            \makecell[{{p{\linewidth}}}]{
                \texttt{\tiny{[GM$|$GM]}}
                \texttt{ДВИГАТЬСЯ: M, (1, 7)} \\
            }
        }
    }
    & & \\ \\

    \theutterance \stepcounter{utterance}  
    & & & \multicolumn{2}{p{0.3\linewidth}}{
        \cellcolor[rgb]{0.9,0.9,0.9}{
            \makecell[{{p{\linewidth}}}]{
                \texttt{\tiny{[GM$|$GM]}}
                \texttt{M moved to (1, 7)} \\
            }
        }
    }
    & & \\ \\

    \theutterance \stepcounter{utterance}  
    & & \multicolumn{4}{p{0.6\linewidth}}{
        \cellcolor[rgb]{0.9,0.9,0.9}{
            \makecell[{{p{\linewidth}}}]{
                \texttt{\tiny{[GM$\rangle$P2]}}
                \texttt{Успешно перемещён объект "M" в (1,7). Ваша обновлённая сетка выглядит так:} \\
\texttt{```} \\
\texttt{ 1234567} \\
\texttt{╔═══════╗} \\
\texttt{║◌◌◌◌◌◌◌║ 1} \\
\texttt{║O◌◌◌◌◌◌║ 2} \\
\texttt{║◌◌◌◌◌┌─╢ 3} \\
\texttt{║S◌◌┌─┤◌║ 4} \\
\texttt{║DA◌│◌├─╢ 5} \\
\texttt{╟───┼─┘◌║ 6} \\
\texttt{║M◌T│P◌◌║ 7} \\
\texttt{╚═══╧═══╝} \\
\texttt{```} \\
\texttt{Сейчас вы играете раунд 10 из максимум 28.} \\
\texttt{Вы вместе накопили 10 из 16 штрафов.} \\
\texttt{Другой игрок переместил объект на своей сетке.} \\
\\ 
\\ 
\texttt{Какая твоя следующая команда?} \\
            }
        }
    }
    & \\ \\

    \theutterance \stepcounter{utterance}  
    & & & \multicolumn{4}{p{0.6\linewidth}}{
        \cellcolor[rgb]{0.9,0.9,0.9}{
            \makecell[{{p{\linewidth}}}]{
                \texttt{\tiny{[GM$\langle$P2]}}
                \texttt{ДВИГАТЬСЯ: T, (1, 3)} \\
            }
        }
    }
     \\ \\

    \theutterance \stepcounter{utterance}  
    & & & \multicolumn{2}{p{0.3\linewidth}}{
        \cellcolor[rgb]{0.9,0.9,0.9}{
            \makecell[{{p{\linewidth}}}]{
                \texttt{\tiny{[GM$|$GM]}}
                \texttt{ДВИГАТЬСЯ: T, (1, 3)} \\
            }
        }
    }
    & & \\ \\

    \theutterance \stepcounter{utterance}  
    & & & \multicolumn{2}{p{0.3\linewidth}}{
        \cellcolor[rgb]{0.9,0.9,0.9}{
            \makecell[{{p{\linewidth}}}]{
                \texttt{\tiny{[GM$|$GM]}}
                \texttt{T moved to (1, 3)} \\
            }
        }
    }
    & & \\ \\

    \theutterance \stepcounter{utterance}  
    & & \multicolumn{4}{p{0.6\linewidth}}{
        \cellcolor[rgb]{0.9,0.9,0.9}{
            \makecell[{{p{\linewidth}}}]{
                \texttt{\tiny{[P1$\langle$GM]}}
                \texttt{Успешно перемещён объект "M" в (1,7). Ваша обновлённая сетка выглядит так:} \\
\texttt{```} \\
\texttt{ 1234567} \\
\texttt{╔═══════╗} \\
\texttt{║P◌◌◌◌◌◌║ 1} \\
\texttt{║O◌◌◌◌◌◌║ 2} \\
\texttt{║T◌◌◌◌┌─╢ 3} \\
\texttt{║S◌◌┌─┤◌║ 4} \\
\texttt{║DA◌│◌├─╢ 5} \\
\texttt{╟───┼─┘◌║ 6} \\
\texttt{║M◌◌│◌◌◌║ 7} \\
\texttt{╚═══╧═══╝} \\
\texttt{```} \\
\texttt{Сейчас вы играете раунд 10 из максимум 28.} \\
\texttt{Вы вместе накопили 10 из 16 штрафов.} \\
\texttt{Другой игрок переместил объект на своей сетке.} \\
\\ 
\\ 
\texttt{Какая твоя следующая команда?} \\
            }
        }
    }
    & \\ \\

    \theutterance \stepcounter{utterance}  
    & \multicolumn{4}{p{0.6\linewidth}}{
        \cellcolor[rgb]{0.9,0.9,0.9}{
            \makecell[{{p{\linewidth}}}]{
                \texttt{\tiny{[P1$\rangle$GM]}}
                \texttt{СКАЗАТЬ: готово?} \\
            }
        }
    }
    & & \\ \\

    \theutterance \stepcounter{utterance}  
    & & & \multicolumn{2}{p{0.3\linewidth}}{
        \cellcolor[rgb]{0.9,0.9,0.9}{
            \makecell[{{p{\linewidth}}}]{
                \texttt{\tiny{[GM$|$GM]}}
                \texttt{СКАЗАТЬ: готово?} \\
            }
        }
    }
    & & \\ \\

    \theutterance \stepcounter{utterance}  
    & & \multicolumn{4}{p{0.6\linewidth}}{
        \cellcolor[rgb]{0.9,0.9,0.9}{
            \makecell[{{p{\linewidth}}}]{
                \texttt{\tiny{[GM$\rangle$P2]}}
                \texttt{Успешно перемещён объект "T" в (1,3). Ваша обновлённая сетка выглядит так:} \\
\texttt{```} \\
\texttt{ 1234567} \\
\texttt{╔═══════╗} \\
\texttt{║◌◌◌◌◌◌◌║ 1} \\
\texttt{║O◌◌◌◌◌◌║ 2} \\
\texttt{║T◌◌◌◌┌─╢ 3} \\
\texttt{║S◌◌┌─┤◌║ 4} \\
\texttt{║DA◌│◌├─╢ 5} \\
\texttt{╟───┼─┘◌║ 6} \\
\texttt{║M◌◌│P◌◌║ 7} \\
\texttt{╚═══╧═══╝} \\
\texttt{```} \\
\texttt{Сейчас вы играете раунд 11 из максимум 28.} \\
\texttt{Вы вместе накопили 10 из 16 штрафов.} \\
\texttt{Другой игрок отправил это сообщение:} \\
\\ 
\\ 
\texttt{"готово?"} \\
\\ 
\\ 
\texttt{Какая твоя следующая команда?} \\
            }
        }
    }
    & \\ \\

    \theutterance \stepcounter{utterance}  
    & & & \multicolumn{4}{p{0.6\linewidth}}{
        \cellcolor[rgb]{0.9,0.9,0.9}{
            \makecell[{{p{\linewidth}}}]{
                \texttt{\tiny{[GM$\langle$P2]}}
                \texttt{СКАЗАТЬ: готово!} \\
            }
        }
    }
     \\ \\

    \theutterance \stepcounter{utterance}  
    & & & \multicolumn{2}{p{0.3\linewidth}}{
        \cellcolor[rgb]{0.9,0.9,0.9}{
            \makecell[{{p{\linewidth}}}]{
                \texttt{\tiny{[GM$|$GM]}}
                \texttt{СКАЗАТЬ: готово!} \\
            }
        }
    }
    & & \\ \\

    \theutterance \stepcounter{utterance}  
    & & & \multicolumn{2}{p{0.3\linewidth}}{
        \cellcolor[rgb]{0.9,0.9,0.9}{
            \makecell[{{p{\linewidth}}}]{
                \texttt{\tiny{[GM$|$GM]}}
                \texttt{true} \\
            }
        }
    }
    & & \\ \\

    \theutterance \stepcounter{utterance}  
    & & & \multicolumn{2}{p{0.3\linewidth}}{
        \cellcolor[rgb]{0.9,0.9,0.9}{
            \makecell[{{p{\linewidth}}}]{
                \texttt{\tiny{[GM$|$GM]}}
                \texttt{Initial states:} \\
\texttt{```} \\
\texttt{Player 1:} \\
\texttt{ 1234567} \\
\texttt{╔═══════╗} \\
\texttt{║◌◌◌◌◌T◌║ 1} \\
\texttt{║P◌◌◌◌◌◌║ 2} \\
\texttt{║◌◌AD◌┌─╢ 3} \\
\texttt{║◌◌◌┌─┤S║ 4} \\
\texttt{║O◌◌│◌├─╢ 5} \\
\texttt{╟───┼─┘◌║ 6} \\
\texttt{║◌M◌│◌◌◌║ 7} \\
\texttt{╚═══╧═══╝} \\
\\ 
\\ 
\texttt{Player 2:} \\
\texttt{ 1234567} \\
\texttt{╔═══════╗} \\
\texttt{║◌◌◌◌◌S◌║ 1} \\
\texttt{║◌◌◌◌O◌◌║ 2} \\
\texttt{║◌◌◌M◌┌─╢ 3} \\
\texttt{║◌◌A┌─┤◌║ 4} \\
\texttt{║◌◌◌│◌├─╢ 5} \\
\texttt{╟───┼─┘◌║ 6} \\
\texttt{║D◌T│P◌◌║ 7} \\
\texttt{╚═══╧═══╝} \\
\\ 
\texttt{```} \\
            }
        }
    }
    & & \\ \\

    \theutterance \stepcounter{utterance}  
    & & & \multicolumn{2}{p{0.3\linewidth}}{
        \cellcolor[rgb]{0.9,0.9,0.9}{
            \makecell[{{p{\linewidth}}}]{
                \texttt{\tiny{[GM$|$GM]}}
                \texttt{End states:} \\
\texttt{```} \\
\texttt{Player 1:} \\
\texttt{ 1234567} \\
\texttt{╔═══════╗} \\
\texttt{║P◌◌◌◌◌◌║ 1} \\
\texttt{║O◌◌◌◌◌◌║ 2} \\
\texttt{║T◌◌◌◌┌─╢ 3} \\
\texttt{║S◌◌┌─┤◌║ 4} \\
\texttt{║DA◌│◌├─╢ 5} \\
\texttt{╟───┼─┘◌║ 6} \\
\texttt{║M◌◌│◌◌◌║ 7} \\
\texttt{╚═══╧═══╝} \\
\\ 
\\ 
\texttt{Player 2:} \\
\texttt{ 1234567} \\
\texttt{╔═══════╗} \\
\texttt{║◌◌◌◌◌◌◌║ 1} \\
\texttt{║O◌◌◌◌◌◌║ 2} \\
\texttt{║T◌◌◌◌┌─╢ 3} \\
\texttt{║S◌◌┌─┤◌║ 4} \\
\texttt{║DA◌│◌├─╢ 5} \\
\texttt{╟───┼─┘◌║ 6} \\
\texttt{║M◌◌│P◌◌║ 7} \\
\texttt{╚═══╧═══╝} \\
\\ 
\texttt{```} \\
            }
        }
    }
    & & \\ \\

    \theutterance \stepcounter{utterance}  
    & & & \multicolumn{2}{p{0.3\linewidth}}{
        \cellcolor[rgb]{0.9,0.9,0.9}{
            \makecell[{{p{\linewidth}}}]{
                \texttt{\tiny{[GM$|$GM]}}
                \texttt{* success: True} \\
\texttt{* lose: False} \\
\texttt{* aborted: False} \\
\texttt{{-}{-}{-}{-}{-}{-}{-}} \\
\texttt{* Shifts: 12.00} \\
\texttt{* Max Shifts: 12.00} \\
\texttt{* Min Shifts: 6.00} \\
\texttt{* End Distance Sum: 7.21} \\
\texttt{* Init Distance Sum: 31.75} \\
\texttt{* Expected Distance Sum: 29.33} \\
\texttt{* Penalties: 10.00} \\
\texttt{* Max Penalties: 16.00} \\
\texttt{* Rounds: 11.00} \\
\texttt{* Max Rounds: 28.00} \\
\texttt{* Object Count: 7.00} \\
            }
        }
    }
    & & \\ \\

\end{supertabular}
}

\end{document}
